%% Generative AI -- Assignment 1 -- technical report
%% IEEE conference format. Build: pdflatex -> bibtex -> pdflatex -> pdflatex
%% (see report/README.md). IEEEtran.cls and IEEEtran.bst are included in this
%% folder so the document compiles anywhere, including Overleaf.
\documentclass[conference]{IEEEtran}

\usepackage[T1]{fontenc}
\usepackage[utf8]{inputenc}
\usepackage{amsmath,amssymb}
\usepackage{graphicx}
\usepackage{booktabs}
\usepackage{array}
\usepackage{multirow}
\usepackage{xcolor}
\usepackage{url}
\usepackage{tikz}
\usetikzlibrary{arrows.meta,positioning,fit}
% Courier's T1 encoding turns "--" into a dash; command-line options such as
% --build and --resume must print as two hyphens.
\usepackage[protrusion=false,expansion=false]{microtype}
\DisableLigatures[-]{encoding=T1,family=tt*}
\usepackage[hidelinks]{hyperref}

% Placeholders for everything that depends on experimental results that do not
% exist yet. Every \todo states exactly which figure, table or number belongs there.
\newcommand{\todo}[1]{\textcolor{red}{[TODO: #1]}}
% Short placeholder for one empty result cell. The note under each results table
% says, with \todo, exactly which output file supplies that table's numbers.
\newcommand{\tbd}{\textcolor{red}{--}}

% A framed red box of the given width and height, shown wherever a figure file does
% not exist yet. \sloppy and the emergency stretch let long file paths break.
\newcommand{\placeholderbox}[3]{% width, height, text
  \fbox{\parbox[c][#2][c]{\dimexpr#1-2\fboxsep-2\fboxrule\relax}{%
    \sloppy\emergencystretch=4em\centering\color{red}\scriptsize #3}}}
\newcommand{\placeholderfig}[1]{\placeholderbox{0.95\columnwidth}{0.16\textheight}{#1}}

% Result figures and tables. Every figure and LaTeX table written by the scripts is
% included through \resultfig / \resulttable: if the file exists in figures/ or
% tables/ it is used, otherwise a red placeholder names the missing file and the
% script output it comes from. Copying the file in (report/collect_outputs.py does
% it for every output) is enough; the text never needs editing.
% \resultfig[placeholder height]{file in figures/}{width}{what it shows and its source}
\newcommand{\resultfig}[4][0.16\textheight]{%
  \IfFileExists{figures/#2}%
    {\includegraphics[width=#3,height=0.42\textheight,keepaspectratio]{figures/#2}}%
    {\placeholderbox{#3}{#1}{\textbf{Missing figure} \path{report/figures/#2}\\[2pt]#4}}}
% \resulttable[maximum width]{file in tables/}{placeholder shown while it is missing}
% A generated file holds a complete table float (booktabs) with its own caption and
% label; only its tabular is kept, so the caption and label written here apply.
% Underscores in generated headers (e.g. EDGE_RATIO) are printed literally, and a
% tabular wider than the available width is scaled down to fit.
\newsavebox{\generatedtablebox}
\let\mathsubscript=_
\begingroup
\catcode`\_=\active
\gdef\textsafeunderscore{\catcode`\_=\active
  \def_{\ifmmode\mathsubscript\else\textunderscore\fi}}
\endgroup
\newcommand{\generatedtabular}[1]{%
  \begingroup
  \renewenvironment{table}[1][]{}{}%
  \renewcommand{\caption}[2][]{}%
  \renewcommand{\label}[1]{}%
  \textsafeunderscore
  \input{#1}%
  \endgroup}
\newcommand{\resulttable}[3][\linewidth]{%
  \IfFileExists{tables/#2}%
    {\sbox{\generatedtablebox}{\generatedtabular{tables/#2}}%
     \ifdim\wd\generatedtablebox>#1\relax
       \resizebox{#1}{!}{\usebox{\generatedtablebox}}%
     \else\usebox{\generatedtablebox}\fi}%
    {#3}}
% Red note under a placeholder table naming the file that will replace it.
\newcommand{\missingtable}[2]{\\[2pt]{\color{red}\scriptsize
  \textbf{Missing table} \path{report/tables/#1}, copied from \path{#2}}}

\newcommand{\repourl}{https://github.com/fsdev87/generative-ai-A1}

\begin{document}

\title{Universal, Hard-Routed and Soft Mixture-of-Experts Image Restoration,\\
and Style-Conditioned Face-to-Sketch Generation}

\author{\IEEEauthorblockN{Mustafa Faraz}
\IEEEauthorblockA{Department of Computer Science\\
National University of Computer and Emerging Sciences (FAST-NUCES), Islamabad\\
mustafafaraz87@gmail.com}}

\maketitle

\begin{abstract}
This report presents four related generative systems built on a single data
pipeline. Clean images of the Oxford-IIIT Pet dataset are corrupted at load time
by one of three degradations---salt-and-pepper noise, Gaussian blur or
rectangular occlusion---or left clean, and three systems of increasing routing
sophistication learn to invert them. The first is a universal convolutional
denoising autoencoder with a genuine compressed bottleneck, which is never told
which corruption it is facing. The second replaces one shared model by a
four-class corruption classifier that routes each image to one of three
specialist autoencoders, with an identity bypass for clean inputs. The third
makes that routing differentiable: a temperature-scaled softmax gate, initialised
from the classifier, assigns continuous weights to the identity branch and the
three specialists, and the whole system is fine-tuned jointly under a loss that
combines reconstruction, structural similarity, corruption classification and a
routing-balance regulariser. The fourth system is a conditional GAN that
translates a facial photograph into a pencil sketch in one of three artist
styles of the FS2K dataset, using a U-Net generator whose decoder is modulated by
a learned style embedding and a PatchGAN discriminator conditioned by projection.
Every model is tuned with Optuna, tracked with Weights \& Biases, exported to
ONNX with a verified parity check, and served through a FastAPI backend behind a
React interface that runs under Docker Compose.
On the 33{,}021 corrupted entries of the official test manifest the deployed
universal autoencoder, which keeps an $8\times8$ bottleneck plus one limited
$32\times32$ skip, reaches 27.93~dB PSNR and 0.879 SSIM without being told the
corruption; an ablation shows that the limited skip restores detail (clean detail
ratio 0.48 to 0.81) while impulse survival \emph{falls}, so it does not leak the
corruption. The corruption classifier reaches 0.998 test accuracy (macro-F1
0.997), hard routing raises the corrupted-input PSNR to 29.12~dB and returns
99\% of clean images exactly, and the soft mixture reaches 29.45~dB, beating hard
routing on 83\% of corrupted entries with nearly one-hot routing. The sketch
generator reaches SSIM 0.46 and $L_1$ 0.110 on the FS2K test set, and swapping the
style raises its $L_1$ error by 38\%.
\end{abstract}

\begin{IEEEkeywords}
image restoration, denoising autoencoder, mixture of experts, conditional
generative adversarial networks, hyperparameter optimisation, ONNX
\end{IEEEkeywords}

\section{Introduction}

Image restoration systems are usually built one degradation at a time. A
denoiser is trained on noise, a deblurring network on blur, an inpainting
network on holes, and the choice between them is left to whoever runs the
pipeline. That arrangement is convenient for research but awkward in
deployment, where the degradation afflicting an incoming image is rarely
announced. Two families of answers have emerged. One trains a single model
across degradations and relies on its capacity to cover all of
them~\cite{zhang2017beyond,li2022allinone,potlapalli2023promptir}; the other
keeps specialised components and adds a mechanism that decides which to
use~\cite{yu2018crafting,yu2022path}.

This assignment asks for both, and for the step that connects them. Using one
fixed data pipeline we build three restoration systems that differ only in how
the knowledge of the corruption enters the computation.

\begin{itemize}
  \item \textbf{Task 1} is a universal denoising autoencoder. It receives a
  corrupted image and nothing else: no label, no severity, no mask. All four
  input conditions---clean, salt-and-pepper, blur, occlusion---must be handled
  by one encoder, one compressed latent representation and one decoder.
  \item \textbf{Task 2} makes the corruption explicit. A convolutional
  classifier predicts which of the four conditions holds, and that discrete
  decision selects one of three specialist autoencoders, or an identity bypass
  when the image is judged clean. The system is evaluated twice: with the true
  label from the test manifest (\emph{oracle routing}, which isolates how good
  the specialists are) and with the classifier's own prediction
  (\emph{predicted routing}, which is the operational system).
  \item \textbf{Task 3} replaces the \texttt{argmax} with a temperature-scaled
  softmax. The gate is the Task~2 classifier, the experts are the Task~2
  specialists, and the output is their weighted sum. Because that sum is
  differentiable, the gate and the experts can be trained jointly on the final
  reconstruction error, with a balance regulariser to keep the gate from
  collapsing onto a single branch.
\end{itemize}

\textbf{Task 4} is a different problem on a different dataset: paired
image-to-image translation with a categorical condition. Given a facial
photograph and one of three sketch styles, a conditional GAN must produce the
sketch that the corresponding artist would have drawn.

Three themes run through the report. The first is \emph{evidence}: every
architectural choice---where the bottleneck sits, whether skip connections are
allowed, which loss, which upsampling operator, which balance regulariser, how
the style condition is injected---is argued from published results and, where
the literature is ambiguous, measured. The second is \emph{honesty about what
the numbers mean}: PSNR is capped for identical images, clean inputs through an
identity bypass are exact rather than merely excellent, FID is biased at small
sample sizes, and the FS2K test set contains only 46 images of one style. Each
of these is stated where the corresponding result is presented rather than
buried. The third is \emph{deployability}: all seven inference models are
exported to ONNX and checked against their PyTorch originals, and the complete
application starts with a single Docker Compose command.

The implementation, including the data preparation, training, search,
evaluation and export scripts, the backend and the frontend, is at
\url{\repourl}. A demonstration video is at \todo{YouTube link to the 5--7 minute
demonstration video}.

\section{Related Work}

\subsection{Denoising autoencoders and restoration networks}

The denoising autoencoder was introduced as a representation learner: an
autoencoder trained to reconstruct a clean input from a stochastically corrupted
one cannot settle for the identity and must capture the structure of the
data~\cite{vincent2008denoising,vincent2010stacked}. The corruption is applied
during training only, which is exactly the runtime scheme used here. The
compressed central layer is what makes the representation
compact~\cite{hinton2006reducing}, and convolutional encoder-decoder stacks make
the idea practical on images~\cite{masci2011stacked,xie2012image}. DnCNN showed
that one network can cover a \emph{range} of noise levels rather than one
model per level~\cite{zhang2017beyond}, the precedent for a single model
covering several corruption families.

Two architectural traditions meet in Task~1. U-Net concatenates every encoder
stage into the mirrored decoder stage~\cite{ronneberger2015unet}, and RED-Net
adds symmetric skip connections to a very deep restoration network, reporting
both faster convergence through better gradient flow and better detail recovery
because the skips ``pass image details from convolutional layers to
de-convolutional layers''~\cite{mao2016image}. That second mechanism is precisely
what a bottleneck is supposed to forbid: a full-resolution skip lets the decoder
reconstruct detail the latent never encoded. Context Encoders take the opposite
route for inpainting, using a narrow bottleneck and no skips, because occluded
content must be synthesised from context rather than copied~\cite{pathak2016context}.
Section~\ref{sec:task1} quantifies both sides of this trade-off for the
bottleneck model rather than assuming them.

\subsection{Routing and all-in-one restoration}

RL-Restore assembles a toolbox of small, differently specialised CNNs and learns
by reinforcement learning which tools to apply and in what
order~\cite{yu2018crafting}; Path-Restore extends the idea to spatially varying
routing, selecting a path per image region with a difficulty-regulated
reward~\cite{yu2022path}. Both report that a set of small specialists plus a
routing policy is more parameter-efficient than one large network---the
hypothesis Task~2 tests in its simplest form, with a single image-level decision
and a supervised classifier instead of a learned policy. The opposite design
removes the router entirely: AirNet infers the degradation with a contrastive
encoder that then guides restoration~\cite{li2022allinone}, and PromptIR encodes
degradation-specific information in lightweight learned prompts that modulate a
single network~\cite{potlapalli2023promptir}. These all-in-one models are the
state of the art for the setting of Task~1; they are trained on far more data
with far larger backbones, and no comparison with them is claimed here.

\subsection{Mixtures of experts}

The mixture of experts originates with Jacobs et al.~\cite{jacobs1991adaptive},
who distinguish a \emph{cooperative} error function, in which the target is
compared with a blend $\sum_i p_i o_i$ of expert outputs, from a
\emph{competitive} one, $\sum_i p_i \lVert d - o_i \rVert^2$, in which each
expert must produce the whole output. They argue that the cooperative form
``does not encourage localization'', because each expert only has to cancel the
residual left by the others, and that it ``tends to lead to solutions in which
many experts are used for each case''. Jordan and Jacobs later gave the
construction its statistical form as a mixture model trained by
EM~\cite{jordan1994hierarchical}; Yuksel et al.\ survey two decades of
variants~\cite{yuksel2012twenty}.

The modern, large-scale form adds sparsity. Shazeer et al.\ introduce noisy
top-$k$ gating and observe the central pathology: ``the gating network tends to
converge to a state where it always produces large weights for the same few
experts'', an imbalance that is self-reinforcing because favoured experts are
trained faster and therefore chosen more often~\cite{shazeer2017outrageously}.
Their remedy is a pair of soft constraints on batch statistics, discussed
quantitatively in Section~\ref{sec:task3}. Switch Transformers simplify the two
losses into one~\cite{fedus2022switch,lepikhin2021gshard}, and Soft MoE removes
the problem structurally: every expert slot receives a weighted average of all
tokens, so the model is ``immune to token dropping and expert unbalance'' and
needs no auxiliary balance loss at all~\cite{puigcerver2024soft}. The system in
Task~3 is dense and soft in the same sense---all four branches always run---but
small enough that the gate, not the capacity, is the object of study. The wider
framing is conditional computation, in which each input traverses only part of
the network~\cite{bengio2015conditional,eigen2013learning}; Eigen et al.\ already
observed the gating-collapse failure mode and addressed it with a hard constraint
early in training, and Bengio et al.\ with a soft constraint on the batchwise
average of each gate---the family to which the regulariser used here belongs.

\subsection{Conditional image-to-image translation}

Conditional GANs supply the condition to both generator and
discriminator~\cite{goodfellow2014generative,mirza2014conditional}. pix2pix is
the canonical paired translation model: a U-Net generator, a $70 \times 70$
PatchGAN discriminator, and an $L_1$ term with weight $\lambda = 100$ added to
the adversarial loss, $L_1$ rather than $L_2$ because it ``encourages less
blurring''~\cite{isola2017image}. Isola et al.\ also report that removing the
conditioning from the discriminator makes the generator collapse to nearly the
same output regardless of input, which is direct evidence for conditioning both
networks.

How a \emph{categorical} condition should enter is a design question with
several published answers. The simplest spatially replicates a label map and
concatenates it to the input~\cite{choi2018stargan}. Conditional instance
normalisation learns a per-condition scale and shift, letting a single network
render many styles~\cite{dumoulin2017learned}, a mechanism generalised as
FiLM~\cite{perez2018film,dumoulin2018featurewise} and applied to language-driven
vision~\cite{devries2017modulating}; AdaIN extends it to arbitrary, unseen
styles~\cite{huang2017arbitrary}. On the discriminator side, the auxiliary
classifier of AC-GAN~\cite{odena2017conditional} competes with the projection
discriminator, which adds an inner product between the embedded condition and
the discriminator's feature vector and is motivated by the role of the condition
in the underlying probabilistic model~\cite{miyato2018cgans}.
Section~\ref{sec:task4} explains, and measures, why this work conditions the
generator by FiLM and the discriminator by projection.

Face photo-to-sketch synthesis itself has a long history, from multiscale Markov
random fields on the CUHK datasets~\cite{wang2009face} to the FS2K benchmark used
here~\cite{fan2022facial}.

\section{Dataset Preparation}
\label{sec:data}

Tasks 1--3 use the Oxford-IIIT Pet dataset~\cite{parkhi2012cats,petsdata}: 7{,}349
photographs of 37 cat and dog breeds. Task~4 uses FS2K~\cite{fan2022facial,fs2kdata},
2{,}104 photograph-sketch pairs in three artist styles. All data preparation is
implemented in \texttt{src/data/} and the resulting splits and manifests are
committed to the repository, so every run of every task sees identical data.

\subsection{Splits}

The official \texttt{trainval} list of the Oxford-IIIT Pet dataset contains
3{,}680 images and the official \texttt{test} list 3{,}669. As the brief
requires, the development portion is divided 80/20 with random seed 42, giving
\textbf{2{,}944 training} and \textbf{736 validation} images; the official test
set of \textbf{3{,}669} images is left untouched until the final evaluation. All
three splits are disjoint and all 37 breeds appear in each: the training split
holds 71--91 images per breed (953 cats, 1{,}991 dogs), the validation split
9--29 per breed (235 cats, 501 dogs) and the test split 88--100 per breed
(1{,}183 cats, 2{,}486 dogs). Breed labels are not used by any model---the clean
image is the reconstruction target---but the per-breed counts confirm that the
random split did not concentrate any breed in one partition.

Every image is converted to RGB and resized to $128 \times 128$ with bicubic
resampling and no crop, then cached once as a \texttt{uint8} array so that
training never re-decodes a JPEG. Pixel values are scaled to $[0,1]$ and stored
as float32 in NCHW order; all restoration models end in a sigmoid.

\subsection{Corruption configuration}

Four input conditions are used, in the fixed order
\texttt{(clean, salt, blur, occlusion)} that indexes the classifier logits, the
mixture weights and every table in this report; index~0 is always the clean or
identity branch. Table~\ref{tab:corruptions} gives the training ranges, from
which severities are sampled uniformly at training time, and the three fixed
severity levels used for testing.

\begin{table*}[t]
\caption{Corruption configuration. Training severities are sampled from the
ranges on every image load; the test severities are fixed by the brief and
stored in the test manifest. Blur strength is the standard deviation in pixels
of the kernel that is actually applied after truncation to $k$ taps.}
\label{tab:corruptions}
\centering
\begin{tabular}{@{}llll@{}}
\toprule
Condition & Training range (sampled per load) & Test levels (low / medium / high) & Test strength \\
\midrule
Clean & image passed through unchanged & --- & --- \\
Salt-and-pepper & $p \sim \mathcal{U}(0.02,\,0.15)$; hit pixels set to 0 or 1 equiprobably & $p = 0.03$ / $0.08$ / $0.15$ & --- \\
Gaussian blur & $k \in \{3,5,7\}$ uniformly, $\sigma \sim \mathcal{U}(0.5,\,2.5)$ & $(3,\,0.7)$ / $(5,\,1.5)$ / $(7,\,2.5)$ & $0.647$ / $1.195$ / $1.765$ px \\
Rectangular occlusion & $n \in \{1,2,3\}$ rectangles covering $\mathcal{U}(0.10,\,0.35)$ of the area & $n=1,2,3$ at $10\%$ / $20\%$ / $35\%$ & --- \\
\bottomrule
\end{tabular}
\end{table*}

Three details of the implementation (\texttt{src/data/corruptions.py}) matter for
the results.

\emph{Salt-and-pepper} noise selects each pixel independently with probability
$p$ and sets all three channels to 0 or 1 with equal probability, so the
corruption is a pure impulse: a hit pixel carries no information about its clean
value.

\emph{Gaussian blur} is separable with reflect padding and matches
\texttt{torchvision}'s \texttt{gaussian\_blur} to within $10^{-5}$, verified by a
unit test. The kernel is truncated to $k$ taps, and this truncation has a
consequence that was discovered during data preparation and changed how
severity is labelled. The standard deviation of the kernel that is actually
applied,
\begin{equation}
  s(k,\sigma) = \Bigl( \textstyle\sum_{i=0}^{k-1} g_i \bigl(i - \tfrac{k-1}{2}\bigr)^2 \Bigr)^{1/2},
  \label{eq:blurstrength}
\end{equation}
in which $g$ denotes the normalised $k$-tap Gaussian kernel,
saturates as $\sigma$ grows, because a wide Gaussian truncated to $k$ taps
flattens towards a box filter whose standard deviation is bounded by
$\sqrt{(k^2-1)/12}$. At $k=3$ that bound is $0.816$~px, so the configuration
$(3,\,2.5)$ has strength $0.805$~px---\emph{milder} than the medium test level
$(5,\,1.5)$ at $1.195$~px, although its $\sigma$ is larger. Ranking blur severity
by $\sigma$ alone is therefore wrong. Severity labels are instead assigned by
binning $s(k,\sigma)$ with edges at the midpoints between the three test levels,
$0.921$ and $1.480$~px. Regenerating the validation manifest under this rule
left every corruption setting bit-identical and corrected \textbf{86 of its 184}
blur severity labels.

\emph{Occlusion} places $n$ non-overlapping black rectangles whose union covers
the target fraction to within one percentage point. The target area is divided
among the rectangles by a Dirichlet draw, each rectangle receives a random
aspect ratio in $[1/2, 2]$, and positions are rejected until no two overlap;
the achieved coverage is recomputed from the union mask and resampled if it
leaves the required range, so the stored coverage is the true one rather than
the request.

A consequence of these definitions is worth stating before any result is
reported: \textbf{all three test ``high'' levels sit at the upper edge of the
training ranges}. Salt-and-pepper high uses $p = 0.15$, exactly the top of the
training range; occlusion high covers $35\%$, exactly the top; and blur high,
$(7,\,2.5)$ at strength $1.765$~px, is the single strongest blur the training
distribution can produce, reachable only at $k=7$ with $\sigma$ near its maximum.
Under uniform sampling of $(k,\sigma)$ only about $15\%$ of training blurs fall
in the ``high'' band at all, and none exceeds the test level. The high severities
are thus the sparsest part of the training distribution, and the weakest results
are expected there---not because the models fail, but because the brief's test
protocol probes the edge of what training covers.

\subsection{Runtime corruption and deterministic manifests}

Training corruptions are generated \emph{at load time}: each access to a training
image draws a fresh condition and a fresh severity, applies the corruption in
NumPy and returns the pair (corrupted input, clean target) together with the
integer label that produced it. Nothing corrupted is ever written to disk, so an
image is seen under a different degradation in every epoch, which acts as a
strong data augmentation and provides the classifier of Task~2 with free,
perfectly accurate labels. A random horizontal flip is applied to the clean
image before corruption.

Validation and test corruptions must instead be identical across runs, models
and machines. They are therefore generated once into JSON \emph{manifests} that
store, for each entry, the corruption type, the severity level, every parameter
($p$; $k$ and $\sigma$; the number of rectangles, achieved coverage and exact
rectangle coordinates) and the random seed. Applying a stored specification is
fully deterministic, so the manifests reproduce the corrupted images exactly
without storing them.

The \textbf{validation manifest} (seed 1001) assigns \emph{one} condition to each
of the 736 validation images, balanced across the four classes: exactly 184
entries each of clean, salt-and-pepper, blur and occlusion. Severities are drawn
from the training ranges, so the validation distribution matches training.
Table~\ref{tab:valmanifest} gives the resulting severity composition. The
distribution across levels is not uniform, and for a well understood reason:
severity is binned on a physical quantity ($p$, covered fraction, blur strength)
whose uniform sampling does not produce uniform bins. For blur the effect is
strongest, since $k=3$ can never reach the medium band at all.

\begin{table}[t]
\caption{Validation manifest (736 entries, seed 1001): one deterministic
corruption per validation image, balanced across the four conditions.}
\label{tab:valmanifest}
\centering
\begin{tabular}{@{}lrrrr@{}}
\toprule
Condition & Low & Medium & High & Total \\
\midrule
Clean & \multicolumn{3}{c}{(no severity)} & 184 \\
Salt-and-pepper & 46 & 97 & 41 & 184 \\
Gaussian blur & 92 & 62 & 30 & 184 \\
Rectangular occlusion & 36 & 91 & 57 & 184 \\
\midrule
Total & & & & 736 \\
\bottomrule
\end{tabular}
\end{table}

The \textbf{test manifest} (seed 2002) is exhaustive instead of sampled: every one
of the 3{,}669 test images appears in ten variants---clean, plus each of the
three corruptions at each of the three fixed severity levels---giving
\textbf{36{,}690 entries}, exactly 3{,}669 per (type, level) group. Because every
image appears under every condition, comparisons across corruptions and
severities are paired. The test set is balanced per (type, level) group but not
per class: the clean class has 3{,}669 entries and each corruption 11{,}007, so
for the Task~2 classifier macro-averaged metrics weight the clean class as much
as each corruption, whereas accuracy (which equals micro-averaged F1 for
single-label classification) is dominated by the corruptions. The achieved
occlusion coverage across the test manifest is tightly controlled: mean $0.1000$,
$0.2000$ and $0.3500$ for the three levels, with extremes of $0.0987$--$0.1016$,
$0.1972$--$0.2028$ and $0.3462$--$0.3544$.

Figure~\ref{fig:corruptiongrid} shows the four conditions at the three fixed test
severities.

\begin{figure*}[t]
\centering
\resultfig{corruption_grid.png}{\textwidth}{Copied from
\path{outputs/corruption_grid.png} (\path{scripts/preview_corruptions.py}): three
test images by ten variants.}
\caption{The four input conditions at the three fixed test severities, applied to
three test images by the deterministic test manifest: clean, salt-and-pepper
($p = 0.03, 0.08, 0.15$), Gaussian blur ($(3,0.7), (5,1.5), (7,2.5)$) and
occlusion (one, two and three rectangles covering 10, 20 and 35\%). Every test
image appears in all ten variants.}
\label{fig:corruptiongrid}
\end{figure*}

\subsection{Class-balanced sampling}

Wherever batches must be balanced---the Task~2 classifier and the Task~3
mixture---a custom batch sampler emits batches of (index, condition) pairs
containing exactly $B/4$ images of each condition, reshuffling the images every
epoch and dropping the final incomplete batch; batch sizes are therefore
multiples of four. Because the condition is drawn by the sampler rather than by
the dataset, perfect balance costs nothing: there is no minority class to
oversample. This is the strongest form of the remedy that Buda et al.\ find
dominant for class imbalance in CNNs, namely oversampling to the point where the
imbalance is eliminated entirely~\cite{buda2018systematic}. It also has a second
consequence that Section~\ref{sec:task3} depends on: in an exactly balanced
batch, a perfectly routing gate assigns mean weight exactly $1/4$ to each
branch, so the routing-balance regulariser is zero at the optimum of the
classification objective instead of fighting it.

subsection{FS2K preparation for Task 4}
abel{sec:fs2kprep}

FS2K provides 2{,}104 photograph-sketch pairs with an official division into
1{,}058 training and 1{,}046 test pairs and a style label in $\{0,1,2\}$ for each
pair~\cite{fan2022facial}. Pairs are matched by image number rather than sorted
position, since the identifiers are not contiguous, and all 2{,}104 pairs were
verified to be one-to-one. Reserving $15\%$ of the official training portion as a
validation split, stratified by style with seed 42, gives \textbf{899 training},
\textbf{159 validation} and \textbf{1{,}046 test} pairs, with styles
$303/298/298$, $54/52/53$ and $619/381/46$ respectively.

Two measured properties of FS2K shape the Task~4 experiment design.

First, \textbf{the official test set is strongly style-imbalanced}: only 46 of
1{,}046 test pairs ($4.4\%$) are style~3, against 619 of style~1. Pooled metrics
are therefore dominated by style~1, and per-style results are mandatory; any
metric requiring a distribution estimate, such as FID, is meaningless on 46
images.

Second, and more subtly, \textbf{style is confounded with photograph source in
the official training split}. FS2K draws its photographs from three sources, and
in the training portion styles 1 and 2 come only from \texttt{photo1}
(CASIA-WebFace~\cite{yi2014learning}, $250 \times 250$ tightly cropped faces)
while style~3 comes only from \texttt{photo2} and \texttt{photo3} (studio
cut-outs and stock photographs). A generator could therefore learn part of
``style~3'' from the appearance of the photograph rather than from the condition
label. The official test set contains combinations never seen in training---179
\texttt{photo3} photographs with style-1 sketches---which makes it a direct test
of whether the style condition generalises to a new photograph domain, exactly
the situation the deployed application faces. Test metrics are therefore
reported per style \emph{and} separately for seen and unseen (source, style)
combinations. A related artefact: the \texttt{photo3} style-3 sketches were drawn
on grey paper (90th-percentile grey level 232 against 250--252 elsewhere), so the
179 \texttt{photo3} style-1 test pairs carry a background offset of roughly ten
grey levels relative to what training taught. The sketches are left unmodified so
that metrics are computed against the official ground truth, and the effect is
reported rather than normalised away.

All FS2K sketches are exactly greyscale ($R = G = B$ in every pixel of all
2{,}104 sketches), so they are stored with one channel. Photographs and sketches
are made EXIF-upright, composited onto white where transparent, centre-square
cropped and resized to $128 \times 128$ with antialiased bicubic resampling;
values are scaled to $[-1,1]$ because the generator ends in a $\tanh$. The centre
crop is used rather than an anisotropic resize or padding because the
non-square sources are exactly the style-3 sources, so either alternative would
add a second shortcut correlated with the style label. Spatial augmentation
(paired horizontal flip, and the pix2pix jitter of upscaling by up to
$286/256$ and cropping back) is applied to the photograph and sketch as one
stacked tensor, so identical geometry is guaranteed by construction.

\section{Task 1: Universal Denoising Autoencoder}
\label{sec:task1}

\subsection{Methodology}

\subsubsection{Architecture}

The universal model (\texttt{src/models/autoencoder.py}, shared with Tasks~2
and~3) is a convolutional autoencoder. A stem of two $3\times3$
convolution-BatchNorm-ReLU layers is followed by $d$ stride-2 blocks that halve
the resolution and double the channel count, capped at eight times the base
width; a $1\times1$ convolution projects the deepest feature map to
$c_\text{lat}$ channels, and a mirrored decoder upsamples back to full
resolution and ends in a sigmoid.

The depth is fixed at $d = 4$, giving an $8 \times 8$ latent. The reason is the
occlusion class: the largest test occlusion is a single rectangle of roughly
$76 \times 76$ pixels, and to synthesise the centre of that hole the network must
see past its border on both sides. Measuring the receptive field by
back-propagating from a single output pixel gives 175 px at $d = 4$---the whole
image---against 87 px at $d = 3$, which barely spans the hole.

The latent is the \textbf{only} path from input to output, asserted by a unit
test that checks the output is invariant to changes in the input once the latent
is held fixed. Its dimension is $c_\text{lat} \times 8 \times 8$ against the
$3 \times 128 \times 128 = 49{,}152$ values of the input, so the compression
ratio ranges from $96\times$ at $c_\text{lat} = 8$ to $12\times$ at
$c_\text{lat} = 64$, where the search range deliberately stops: wider latents
would keep improving PSNR, since more capacity always helps a pixel metric, but
the brief requires a \emph{meaningful} bottleneck.

Upsampling uses bilinear resize followed by convolution rather than transposed
convolution. Transposed convolution suffers ``uneven overlap'' whenever the
kernel size is not divisible by the stride, and the unevenness compounds across
the two spatial axes into the familiar checkerboard pattern;
resize-convolution avoids it because it is implicitly weight-tied in a way that
discourages high-frequency artefacts~\cite{odena2016deconvolution}. (An
alternative fix initialises sub-pixel convolution to be equivalent to
nearest-neighbour resize-convolution~\cite{aitken2017checkerboard,shi2016real};
since the decoder here convolves after every upsampling step at every scale, the
simpler resize-convolution was used.)

\subsubsection{Skip connections}

Unrestricted skips, as in U-Net~\cite{ronneberger2015unet}, would let the decoder
copy the input around the bottleneck, which the brief rules out; and for this task
a full-resolution skip is actively harmful, because here the input \emph{is} the
corruption---it hands salt-and-pepper impulses and black occlusion rectangles
straight to the decoder. A first model without any skip, however, plateaued at an
output quality that was nearly independent of the input and lost about half of the
fine detail even on clean images (the $8\times8$ latent is the only path for
texture). Since skips are standard in restoration and RED-Net attributes real gains
to them~\cite{mao2016image}, the final Task~1 models were trained as an ablation
that measures their effect rather than assuming it: the selected configuration was
trained \emph{three} times with identical hyperparameters and epoch budget---no skip
(\texttt{udae}), one skip at $16\times16$ (\texttt{udae\_skip16}) and one at
$32\times32$ (\texttt{udae\_skip32})---and all three were evaluated on the full test
manifest. Skips are only ever taken from deep, low-resolution encoder levels, never
from full resolution, and the latent remains the only path at the coarsest scale, so
every variant keeps a genuine compressed bottleneck. Even the $16\times16$ skip carries
$256 \times 16 \times 16 = 65{,}536$ values at base width 32 (98{,}304 at the
selected width 48), more than the input itself, which makes the point that a
``limited'' skip is limited in \emph{resolution}, not in capacity: deep features
have passed through many nonlinearities and carry context, whereas a
$128\times128$ skip is two convolutions from raw pixels. The variants are
implemented\IfFileExists{tables/task1_skip_ablation.tex}{ and compared in
Section~\ref{sec:task1skip}.}{, but the GPU budget did not allow training them
(Section~\ref{sec:task1skip}).}

Two diagnostics quantify the trade-off directly. \emph{Impulse survival}, on
salt-and-pepper entries, is the fraction of impulse-hit pixels whose output
value is still closer to the impulse than to the clean value; it is 1 for the
identity and 0 for perfect restoration. \emph{Detail ratio}, on clean entries, is
the mean absolute Laplacian of the output luminance divided by that of the
target; below 1 means the output is smoother than the clean image, i.e.\ detail
was lost in the bottleneck. Together they separate the two competing effects:
skips should raise the detail ratio towards 1 and, if they leak, also raise
impulse survival.

\subsubsection{Loss}

The loss is the one prescribed by the brief,
\begin{equation}
  \mathcal{L}_\text{UDAE} = \alpha\,\mathcal{L}_{L_1}(x,\hat{x}) + (1-\alpha)\bigl(1 - \mathrm{SSIM}(x,\hat{x})\bigr),
  \label{eq:udaeloss}
\end{equation}
with SSIM computed as specified by Wang et al.~\cite{wang2004image}: an
$11\times11$ circularly symmetric Gaussian window of standard deviation 1.5
samples, $K_1 = 0.01$, $K_2 = 0.03$, over the valid region and averaged across
channels. The $L_1$ term keeps colours and brightness correct; the SSIM term
rewards local structure and contrast. The motivation for mixing them is Zhao et
al.'s systematic comparison, which finds that $L_2$ produces ``splotchy''
artefacts in flat regions because it tolerates small errors regardless of local
structure, that $L_1$ converges to a better optimum, and that SSIM-type losses
are insensitive to uniform biases and therefore permit colour and brightness
shifts if used alone~\cite{zhao2017loss}. Their recommended mixture uses
$\alpha = 0.84$ \emph{on the MS-SSIM term}, chosen so the two contributions are
roughly balanced; in the brief's parameterisation, where $\alpha$ weights $L_1$,
that corresponds to $\alpha \approx 0.16$. The two conventions must not be
confused, and $\alpha$ is tuned rather than assumed, over a range that contains
both the brief's $0.8$ and Zhao et al.'s balance point. Single-scale SSIM is used
rather than MS-SSIM~\cite{wang2003multiscale} because at $128\times128$ a
five-level pyramid would reduce the coarsest scale to $8\times8$, below the
$11\times11$ window. A perceptual, feature-space
loss~\cite{johnson2016perceptual} was also considered and rejected: it optimises
a quantity different from the PSNR and SSIM that are reported, and adds a VGG
forward pass to every training step on a time-limited GPU.

\subsubsection{Training procedure}

Optimisation uses AdamW~\cite{loshchilov2019decoupled}, which applies weight
decay to the weights directly rather than through the adaptive gradient, with a
fixed decay of $10^{-4}$: overfitting is not the principal risk when every epoch
sees freshly drawn corruptions. The schedule is one epoch of linear warm-up,
which avoids large early Adam steps while the second-moment estimates and the
BatchNorm statistics are still noisy~\cite{goyal2017accurate}, followed by
per-step cosine decay to $1\%$ of the peak rate~\cite{loshchilov2017sgdr}. Mixed
precision (fp16 autocast with gradient scaling~\cite{micikevicius2018mixed}) is
used on the GPU, since the T4 has fp16 tensor cores but no bf16 support, with
SSIM always evaluated in float32 because its variance terms lose all precision
in fp16. Gradients are clipped at norm 1.0 as a safety net, and a non-finite
epoch loss aborts the run (or prunes the Optuna trial) rather than silently
producing NaNs. Spatial dropout~\cite{tompson2015efficient,srivastava2014dropout}
after each convolution block is searched but expected to be of little use, since
runtime corruption is already a strong regulariser.

Model selection and the Optuna objective use a validation score independent of
the tuned loss weights,
\begin{equation}
  \text{score} = 0.5\,\mathrm{SSIM} + 0.5\,\frac{\mathrm{PSNR}}{40},
  \label{eq:score}
\end{equation}
so that trials with different $\alpha$ remain comparable---otherwise the search
would partly be optimising its own yardstick. Both components are reported
separately as well, since PSNR and SSIM measure different things and can
disagree~\cite{hore2010image}.

\subsubsection{Optuna search design}

The study \texttt{task1\_udae} maximises~(\ref{eq:score}) on the validation
manifest with a TPE sampler~\cite{bergstra2011algorithms,akiba2019optuna} seeded
at 42 after ten random start-up trials, and a median
pruner~\cite{golovin2017google} that is inactive for the first five trials and
the first two epochs of each trial. Trial~0 is the brief's own starting
configuration, so the report can state how much the search gained over it. The
search space is given in Table~\ref{tab:task1search}.

\begin{table}[t]
\caption{Optuna search space of Task 1 (study \texttt{task1\_udae}, TPE sampler,
median pruner, 15 epochs per trial), the brief's starting point (trial~0) and the
configuration selected by the search (trial~32).}
\label{tab:task1search}
\centering
\setlength{\tabcolsep}{4pt}
\begin{tabular}{@{}lllll@{}}
\toprule
Parameter & Range & Scale & Trial 0 & Selected \\
\midrule
Learning rate & $10^{-4}$--$3\times10^{-3}$ & log & $10^{-3}$ & $8.28\times10^{-4}$ \\
Batch size & $\{16, 32, 64\}$ & cat. & 32 & \textbf{16} (edge) \\
Latent channels & $\{8, 16, 32, 64\}$ & cat. & 32 & \textbf{64} (edge) \\
Base channels & $\{16, 32, 48, 64\}$ & cat. & 32 & 48 \\
Dropout & $0.0$--$0.3$ & uniform & 0 & 0.016 \\
$\alpha$ ($L_1$ weight) & $0.1$--$0.95$ & uniform & 0.8 & \textbf{0.285} \\
\midrule
\multicolumn{3}{@{}l}{Validation score~(\ref{eq:score}) after 15 epochs} & 0.5940 & \textbf{0.6263} \\
\bottomrule
\end{tabular}
\end{table}

Trials run 15 epochs, a quarter of the 60-epoch final run, and the cosine
schedule anneals completely within each trial so that configurations are
compared as fully annealed models rather than mid-schedule. The known
limitation of this low-fidelity protocol is that settings which pay off only
late---a lower learning rate, or dropout---can be under-rated. The ranges, in
the order of Table~\ref{tab:task1search}: the learning rate brackets Adam's usual
$10^{-3}$; the batch sizes give 184, 92 or 46 optimiser steps per epoch on 2{,}944
images; the latent widths give $96\times$ to $12\times$ compression; the base widths
give 1.0 to 16.5\,M parameters; the dropout range stays low because runtime
corruption already regularises; and the $\alpha$ range spans SSIM- to
$L_1$-dominated losses.

\subsubsection{Search results}

The study was launched for 40 trials with a 90-minute cap, and the cap ended it:
33 trials were run (15 completed, 18 stopped by the median pruner) in 89.6
minutes of trial time on the T4. The best trial was the \emph{last} one, trial~32,
with a validation score of 0.6263 against 0.5940 for the brief's starting
configuration (Fig.~\ref{fig:task1optuna}, left), so the search was still
improving when it stopped and more trials could plausibly have helped. Four
observations explain the selected configuration.

\emph{Loss weighting.} Optuna moved $\alpha$ from the brief's 0.8 to 0.285, putting
most of the weight on the $(1-\mathrm{SSIM})$ term, and $\alpha$ is the most
important parameter in Optuna's importance estimate over the completed trials
(0.44, Fig.~\ref{fig:task1optuna}, right). This agrees with Zhao et al.\ once their
parameterisation is read correctly: their $\alpha = 0.84$ weights the MS-SSIM
term~\cite{zhao2017loss}, which corresponds to $\alpha \approx 0.16$ here, not to
the brief's 0.8 on $L_1$. One caveat applies: the selection score~(\ref{eq:score})
is itself half SSIM, which favours SSIM-heavy losses.

\emph{Bottleneck at the edge of the range.} The best trial used the widest latent
in the range, $64 \times 8 \times 8 = 4{,}096$ values ($12\times$ compression). A
wider bottleneck passes more detail, so reconstruction quality rises with it;
the margin is small here (the best 32-channel trial, trial~24, scored 0.6225) and
the latent width has a low importance (0.02), but the direction is the expected
one. The range was capped deliberately: much wider latents approach an identity
mapping, which the brief rules out by requiring a meaningful bottleneck, so the
edge result reflects a quality-versus-compression trade-off rather than an
oversight of the search.

\emph{Short-budget bias.} The smallest batch size won, and TPE converged on it
(11 of the last 12 trials used batch~16). With a fixed 15-epoch budget, batch~16
gives four times as many optimiser steps as batch~64 (184 against 46 per epoch),
so short trials favour configurations that learn fast; the final run, four
times longer, mitigates but does not remove this bias.

\emph{Dropout.} Dropout has the second-highest importance (0.40), but in the
negative direction: the five best trials all used a rate of at most 0.065 and the
selected value is 0.016, consistent with the expectation that runtime
corruption already regularises the model.

\begin{figure*}[t]
\centering
\resultfig{task1_optuna_optimization_history.png}{0.42\textwidth}{From
\path{outputs/task1/optuna/task1_udae/optimization_history.png}.}\hfil
\resultfig{task1_optuna_param_importances.png}{0.42\textwidth}{From
\path{outputs/task1/optuna/task1_udae/param_importances.png}.}
\caption{Optuna study \texttt{task1\_udae}. Left: validation score of every
completed trial and the running best; pruned trials have no final value and are
not shown. Right: parameter importances estimated from the completed trials.}
\label{fig:task1optuna}
\end{figure*}

\subsubsection{Baselines}

Two reference points accompany every Task~1 result. The \emph{identity} baseline
is the corrupted input itself and measures how damaging each corruption is. The
\emph{oracle classical} baseline applies, to each entry, the textbook method for
its corruption \emph{using information the network never receives}: a $3\times3$
median filter for salt-and-pepper~\cite{hwang1995adaptive,chan2005salt}; unsharp
masking with the true kernel $(k,\sigma)$,
$\hat{x} = y + (y - h \ast y)$, which is one step of Van Cittert-style iterative
deconvolution~\cite{richardson1972bayesian,lucy1974iterative}; and Telea
fast-marching inpainting of the true rectangle
mask~\cite{telea2004image,bertalmio2000image,bertalmio2001navier}. Clean inputs
are returned unchanged, so the oracle is exact there. Because the oracle is given
the corruption type and its exact parameters, beating it is a strong result and
losing to it on one corruption is an informative finding rather than a failure
to conceal.

\subsection{Results}

\subsubsection{Final training}
The configuration selected by the search (\texttt{configs/task1.yaml}) was
trained three times on Kaggle for 60 epochs each, with identical hyperparameters:
without skips (\texttt{udae}, 9.31\,M parameters), with one skip at $16\times16$
(\texttt{udae\_skip16}, 10.64\,M) and with one skip at $32\times32$
(\texttt{udae\_skip32}, 9.64\,M). All three keep the $64\times8\times8$ latent
(4{,}096 values, $12\times$ compression). The best checkpoints are from epochs 56,
59 and 59, so the cosine schedule was still improving the models in its final
low-learning-rate epochs and early stopping never triggered. The deployed model
was chosen by \emph{validation} score~(\ref{eq:score}), never by test results:
0.6946 for \texttt{udae}, 0.7586 for \texttt{udae\_skip16} and 0.8187 for
\texttt{udae\_skip32}, which is therefore the model served by the application and
reported below. Its curves (Fig.~\ref{fig:task1curves}) show training and
validation losses tracking each other closely, as expected when every epoch draws
fresh corruptions, and occlusion as the lowest validation curve throughout.

\begin{figure*}[t]
\centering
\resultfig{task1_training_curves.png}{0.92\textwidth}{From
\path{outputs/task1/figures/training_curves.png}.}
\caption{Task 1 training: training and validation loss (left), and validation
PSNR (centre) and SSIM (right) per input condition on the validation manifest.
The dotted line marks the best epoch (59).}
\label{fig:task1curves}
\end{figure*}

\subsubsection{Test results}
Tables~\ref{tab:task1bytype} and~\ref{tab:task1bylevel} give the test results of
\texttt{udae\_skip32} on the 36{,}690 entries. Pooled over the 33{,}021 corrupted
entries it reaches \textbf{27.93~dB / 0.879}, against 19.26~dB / 0.635 for the
corrupted input and 27.53~dB / 0.861 for the oracle classical baselines, which are
told the corruption, its parameters and the occlusion mask. The blind model thus
beats the informed baselines on average: clearly on salt-and-pepper (31.49 against
29.23~dB), slightly in SSIM on blur (0.873 against 0.855, though 0.9~dB lower in
PSNR), and it is close on occlusion (23.15 against 23.34~dB), where content behind
the mask must be synthesised. Quality falls with severity for blur (32.06, 29.18,
26.22~dB) and occlusion (25.89, 23.05, 20.51~dB) but hardly for salt-and-pepper
(32.01 to 30.85~dB); the high levels sit at the edge of the training ranges
(Section~\ref{sec:data}). Clean inputs come back at 32.21~dB / 0.956: still not
unchanged---everything except the $32\times32$ features passes the bottleneck---and
for the mildest blur the output (32.06~dB) is slightly \emph{below} the corrupted
input (32.34~dB). That residual cost of universality is what the identity branch
of Tasks~2 and~3 removes. Salt-and-pepper impulse survival is 0.0009 and the clean
detail ratio 0.81 (Table~\ref{tab:task1skip}).

\begin{table*}[t]
\caption{Task 1 test results per input condition (36{,}690 entries; 3{,}669 clean
and 11{,}007 per corruption): PSNR (dB) and SSIM of the corrupted input itself, of
the oracle classical baselines that are told the corruption and its parameters,
and of the UDAE, which is told nothing. Both baselines return clean inputs
unchanged (exact, $\infty$); the pooled ``corrupted'' row covers the three
corruptions only.}
\label{tab:task1bytype}
\centering
\footnotesize
\resulttable{task1_comparison_by_type.tex}{%
\begin{tabular}{@{}lcccccc@{}}
\toprule
Type & \multicolumn{2}{c}{Corrupted input} & \multicolumn{2}{c}{Oracle classical} & \multicolumn{2}{c}{UDAE} \\
\midrule
clean & \tbd & \tbd & \tbd & \tbd & \tbd & \tbd \\
\bottomrule
\end{tabular}%
\missingtable{task1_comparison_by_type.tex}{outputs/task1/tables/comparison_by_type.tex}}
\end{table*}

\begin{table*}[t]
\caption{Task 1 test results per input condition and severity level (3{,}669 entries
per row), with the same two reference points as Table~\ref{tab:task1bytype}.}
\label{tab:task1bylevel}
\centering
\footnotesize
\resulttable{task1_comparison_by_type_level.tex}{%
\begin{tabular}{@{}llcccccc@{}}
\toprule
Type & Severity & \multicolumn{2}{c}{Corrupted input} & \multicolumn{2}{c}{Oracle classical} & \multicolumn{2}{c}{UDAE} \\
\midrule
salt & low & \tbd & \tbd & \tbd & \tbd & \tbd & \tbd \\
\bottomrule
\end{tabular}%
\missingtable{task1_comparison_by_type_level.tex}{outputs/task1/tables/comparison_by_type_level.tex}}
\end{table*}

\begin{figure*}[t]
\centering
\resultfig[0.30\textheight]{task1_representative_examples_1.png}{0.48\textwidth}{From
\path{outputs/task1/figures/representative_examples_1.png}.}\hfill
\resultfig[0.30\textheight]{task1_representative_examples_2.png}{0.48\textwidth}{From
\path{outputs/task1/figures/representative_examples_2.png}.}
\caption{Representative Task 1 test restorations, rows of target, corrupted input,
UDAE output and absolute error map, labelled with output PSNR (input PSNR) and
SSIM. Left: three clean entries and salt-and-pepper low/medium/high; right: blur
and occlusion at low/medium/high. Each entry is the one closest to its group's
median SSIM, so the examples are typical rather than chosen by hand.}
\label{fig:task1examples}
\end{figure*}

\begin{figure*}[t]
\centering
\resultfig{task1_severity_curves.png}{0.74\textwidth}{From
\path{outputs/task1/figures/severity_curves.png}.}
\caption{Task 1 restoration quality against severity: PSNR (top) and SSIM (bottom)
of the UDAE, the oracle classical baselines and the corrupted input for each
corruption; the dotted line is the UDAE's score on clean inputs. The UDAE curve is
almost flat at the clean-input level for salt-and-pepper and mild blur---the
${\approx}25$~dB ceiling.}
\label{fig:task1severity}
\end{figure*}

\subsubsection{Skip-connection ablation}
Table~\ref{tab:task1skip} and Fig.~\ref{fig:task1skipexamples} answer the question
the ablation was designed for. Moving one skip from nowhere to $16\times16$ to
$32\times32$ raises corrupted-input quality from 23.48~dB / 0.779 to 25.47~dB /
0.841 to 27.93~dB / 0.879 and clean-input quality from 24.65 to 27.67 to
32.21~dB, and the clean detail ratio from 0.48 to 0.66 to 0.81: the skip passes
spatial detail that the $8\times8$ latent cannot hold. The predicted price---the
corruption leaking through---does \emph{not} appear: impulse survival falls from
0.0024 to 0.0014 to 0.0009. A $32\times32$ feature map has passed through two
stride-2 blocks and carries features, not raw pixels, so an isolated impulse or
the edge of a black rectangle is not copied to the output; the model restores
rather than copies, and the $8\times8$ latent remains the only path for context.
The skips carry two and four times as many values as the input, which is why the
deployed model is described honestly as a bottleneck autoencoder \emph{with one
limited skip}; the pure-bottleneck variant is kept as the reference that shows
what the bottleneck alone achieves. The $128\times128$ skip was not trained: it is
two convolutions from the raw pixels, the case in which leaking is expected.

\begin{figure}[t]
\centering
\resultfig[0.2\textheight]{task1_skip_ablation_examples.png}{\columnwidth}{From
\path{outputs/task1/figures/skip_ablation_examples.png}.}
\caption{The same test entries restored by the three Task~1 variants (no skip,
$16\times16$ skip, $32\times32$ skip).}
\label{fig:task1skipexamples}
\end{figure}
\label{sec:task1skip}

\begin{table}[t]
\caption{Skip-connection ablation: one limited skip added at the stated
resolution, same hyperparameters and epoch budget as the main model.}
\label{tab:task1skip}
\centering
\footnotesize
\resulttable{task1_skip_ablation.tex}{}
\end{table}

\subsection{Failure analysis}
Figure~\ref{fig:task1failures} shows the failures, chosen by rank rather than by
hand. \emph{Black backgrounds read as occlusion:} the worst clean (11.3~dB /
0.492), salt-and-pepper (12.5~dB) and occlusion (10.3~dB) outputs are the same cat
photographed against a uniformly black background, which the model fills with a
grey-green haze exactly as it fills occlusion rectangles; in training, large black
regions are almost always masks, and a blind model cannot tell black content from
a black mask (a limitation of the single-model design and the training
distribution). \emph{Unrecoverable high frequencies:} the worst high-blur output
(22.2~dB from 22.5~dB, SSIM 0.359) is a dog in front of a finely striped blind,
whose stripes the strongest blur erased and the model cannot reinvent.
\emph{No way to leave an image alone:} an occlusion rectangle that falls on a
black background leaves the input practically unchanged (87.3~dB), yet the model
returns 34.9~dB, and a mildly blurred grey cat on a saturated purple backdrop
drops from 37.0 to 26.4~dB. Both are the price of universality, which the identity
branch of Tasks~2 and~3 removes.

\begin{figure}[t]
\centering
\resultfig[0.30\textheight]{task1_failure_cases.png}{\columnwidth}{From
\path{outputs/task1/figures/failure_cases.png}.}
\caption{Task 1 failure cases, selected automatically: the lowest-SSIM output for
each input condition (rows 1--4) and the two corrupted entries on which the UDAE
loses most PSNR relative to its own input (rows 5--6). Labels give output PSNR
(input PSNR) and SSIM.}
\label{fig:task1failures}
\end{figure}

\section{Task 2: Corruption Classification and Hard Routing}
\label{sec:task2}

\subsection{Methodology}

\subsubsection{Classifier}

The corruption classifier (\texttt{src/models/classifier.py}) maps an RGB image
to four logits. It is a conventional convolutional
classifier~\cite{lecun1998gradient,he2016deep} with batch
normalisation~\cite{ioffe2015batch}: each stage applies one or two $3\times3$
convolution-BatchNorm-ReLU layers followed by $2\times2$ max pooling, and two
design choices follow directly from what the cues look like. The first stage runs
at full $128\times128$ resolution, because the evidence is high-frequency: a
few isolated impulse pixels, or the \emph{absence} of high frequencies under mild
blur, would both be destroyed by early downsampling. The head concatenates global
average and global max pooling, because the cues are sparse---a handful of bright
pixels, or one black rectangle---and average pooling alone dilutes them.

Training uses plain multiclass cross-entropy, as the brief requires, with no
label smoothing: the probabilities are displayed in the application and are used
to initialise the Task~3 gate, so they should be the model's own estimates.
Batches are exactly class-balanced by the sampler of Section~\ref{sec:data}, and
the labels come free from the corruption pipeline. AdamW is used with weight
decay applied only to convolution and linear weights, not to BatchNorm
parameters or biases, where decaying them does not regularise what the
normalised layer computes~\cite{he2019bag}.

Model selection, early stopping and the Optuna objective all use validation
\textbf{macro-F1} rather than loss or accuracy. Routing uses the \texttt{argmax},
so the decision metric is what matters; macro-averaging weights the four classes
equally and penalises a class that is over-predicted, which accuracy on a
balanced set would hide~\cite{sokolova2009systematic}. Validation cross-entropy
is logged but not used for selection, because it can rise through growing
over-confidence while accuracy still improves~\cite{guo2017calibration}.

\subsubsection{Specialists}

Three specialist autoencoders are trained, one per corruption, each seeing
\emph{only} its own corruption with a fresh severity on every load and the clean
image as target. The three runs are fully independent---different seeds, separate
optimisers, separate checkpoints---so their parameters are independently trained
as the brief requires. They share the architecture of Section~\ref{sec:task1}.
A first set of specialists with a pure bottleneck produced over-smoothed outputs and
blurry occlusion fills, so the final specialists add \emph{one} limited skip at
$32\times32$ (\texttt{configs/task2\_specialists\_skip.yaml}), one level above
their $16\times16$ latent---the counterpart of Task~1's $16\times16$ skip on its
$8\times8$ latent. No full-resolution skip is used, so a specialist cannot pass its
corruption's pixel-level signature straight to the output; the latent remains the
only path at the coarsest scale.

To keep the experiment computationally feasible the brief permits one shared
architecture search. Each trial trains all three specialists with the trial's
hyperparameters, \emph{interleaved epoch by epoch}, and reports the running mean
of the three validation scores to the pruner. Interleaving matters: training them
sequentially would show the pruner only the salt-and-pepper score for the first
third of each trial, so a configuration that suits impulse noise but fails on
occlusion would survive longer than it deserves. After the study the three
specialists are retrained independently on the full schedule.

Unlike Task~1, the specialist search may also choose the latent \emph{resolution}:
a $16\times16$ latent (encoder depth 3) keeps more spatial layout than an $8\times8$
latent with the same number of values. The six presets, from $8\times8\times32$ to
$16\times16\times64$, hold 2{,}048 to 16{,}384 values ($24\times$ to $3\times$
compression), so every option is still a genuine bottleneck.

\subsubsection{Hard routing}

At inference the classifier produces $p = \mathrm{softmax}(C(\tilde{x}))$ and the
route is $r = \arg\max_k p_k$. The restorer starts from a copy of the input and
overwrites only the entries with $r > 0$ using the corresponding specialist,
which has three consequences worth stating: a clean-routed input is returned
\emph{bit-identical} to the input through the identity bypass, never
unnecessarily processed by a restoration expert; each expert runs at most once
per batch on its own sub-batch rather than once per sample; and all models run in
evaluation mode, so BatchNorm uses running statistics and a sample's output does
not depend on the rest of its batch.

The system is evaluated in both modes the brief prescribes. \emph{Oracle
routing} takes $r$ from the manifest label and isolates the restoration ability
of the specialists; \emph{predicted routing} uses the classifier and measures the
complete operational system. Correctly routed entries give identical metrics in
both modes, which is asserted by a test.

\subsubsection{Optuna search design}

Two studies are run (Table~\ref{tab:task2search}): \texttt{task2\_classifier}
(30 trials of 12 epochs, objective validation macro-F1, 45-minute cap) and
\texttt{task2\_specialists} (20 trials of 6 epochs per specialist, objective the
mean of the three best validation restoration scores, 80-minute cap). Both use
the TPE sampler seeded at 42 with a median pruner. The classifier's channel
presets are \texttt{16-32-64-128}, \texttt{32-64-128-256},
\texttt{48-96-192-384} and the five-stage \texttt{32-64-128-256-256}, which trade
width against depth at 0.21--1.89 GMAC per image; the batch sizes are multiples
of four so that every batch is exactly balanced.

\begin{table}[t]
\caption{Optuna search spaces for Task 2 and the selected configurations (best
trials 1 and 10), committed as \texttt{configs/task2\_classifier.yaml} and
\texttt{configs/task2\_specialists.yaml}.}
\label{tab:task2search}
\centering
\setlength{\tabcolsep}{4pt}
\begin{tabular}{@{}lll@{}}
\toprule
Parameter & Range & Selected \\
\midrule
\multicolumn{3}{@{}l}{\emph{Classifier} (\texttt{task2\_classifier}, 30 trials $\times$ 12 epochs)} \\
Learning rate & $10^{-4}$--$3\times10^{-3}$ (log) & $1.70\times10^{-3}$ \\
Batch size & $\{32, 64, 128\}$ & 32 \\
Channels & four presets & \texttt{32-64-128-256} \\
Convs per stage & $\{1, 2\}$ & 1 \\
Dropout (head) & $0.0$--$0.5$ & 0.146 \\
Weight decay & $10^{-5}$--$10^{-1}$ (log) & $2.9\times10^{-4}$ \\
Validation macro-F1 & & 0.9959 \\
\midrule
\multicolumn{3}{@{}l}{\emph{Specialists} (\texttt{task2\_specialists}, 20 trials $\times$ 6 epochs)} \\
Learning rate & $10^{-4}$--$3\times10^{-3}$ (log) & $2.37\times10^{-3}$ \\
Batch size & $\{16, 32, 64\}$ & \textbf{16} (edge) \\
Base channels & $\{16, 32, 48\}$ & 48 \\
Latent shape & six presets & $16\times16\times32$ ($6\times$) \\
$\alpha$ ($L_1$ weight) & $0.4$--$0.95$ & \textbf{0.516} \\
Mean validation score & & 0.6819 \\
\bottomrule
\end{tabular}
\end{table}

\subsubsection{Search results}

\emph{Classifier.} All 30 trials ran, and took only 15.6 minutes of trial time,
well inside the cap: the median pruner stopped 25 of them, mostly after their
fourth epoch, and only five completed. The best trial was trial~1, with
validation macro-F1 0.9959 (lr $1.70\times10^{-3}$, batch 32, channels
\texttt{32-64-128-256}, one convolution per stage, dropout 0.146, weight decay
$2.9\times10^{-4}$). The pruning rate is itself a finding: the objective is nearly
saturated. Even pruned trials had reached macro-F1 of about 0.98 within three or
four epochs (trial~29: 0.9783), and the five completed trials span only
0.992--0.996 (Fig.~\ref{fig:task2optuna}, left). Detecting the corruption type is
easy for a small CNN on this data, so the search mostly selected an
\emph{inexpensive} configuration---the narrower four-stage network with one
convolution per stage, 0.39\,M parameters, which is also the cheaper gate for
Task~3. With five completed trials so close together, the importance estimate of
this study is not informative and is not used.

\emph{Specialists.} Exactly 20 trials finished, 12 completed and 8 pruned, none
failed, in 61.5 minutes of trial time. The best was trial~10, with a mean
validation score of 0.6819 (salt-and-pepper 0.7006, blur 0.7295, occlusion
0.6155), and its configuration is shared by the three specialists: lr
$2.37\times10^{-3}$, batch 16, base width 48, latent $16 \times 16 \times 32 = 8{,}192$
values ($6\times$ compression) and $\alpha = 0.516$; each specialist has 3.97\,M
parameters. The learning rate dominates the importance estimate (0.76), followed
by $\alpha$ (0.20; Fig.~\ref{fig:task2optuna}, right). The study also gives a free
estimate of run-to-run noise. The first session was stopped after about eleven
minutes, while the checkpoint incident described in the AI-use appendix was being
fixed; trials 0--2 had been saved and the resumed study continued from trial~3.
Because the TPE sampler was re-created with the same seed, it proposed the
parameters of trials 0--2 again as trials 3--5, and the repeated runs differ by
only $2$--$4\times10^{-4}$ in score (0.6046 against 0.6042, 0.6045 against
0.6047). Score differences of that size between trials are therefore noise:
trials 10, 12 and 13 (0.6819, 0.6810, 0.6806), which share every categorical
choice and differ only slightly in learning rate and $\alpha$, are effectively
tied, so the selected configuration is a robust region rather than a single lucky
trial. Three of the 20 trials were spent on these replays, a cost of the resume
logic worth fixing.

Three observations carry over from Task~1. \emph{Loss weighting:} $\alpha$ again
settled well below the brief's 0.8 (0.516; Task~1 chose 0.285), so the searches
consistently moved weight towards the SSIM term. \emph{Short-budget bias:} the
smallest batch size won again. \emph{Speed of learning:} after only six epochs, the
best trial's blur specialist already slightly exceeded the fully trained
universal model on the same 184 validation blur entries (score 0.7295 against
0.7245 after 60 epochs).

\textbf{A confound for the cross-task comparison.} The specialists selected a
latent twice as large as Task~1's: 8{,}192 values ($6\times$ compression) against
4{,}096 ($12\times$). Both are genuine bottlenecks, but part of any advantage of
hard routing over the universal model may come from the larger bottleneck
rather than from specialisation alone; the Task~1 search range, capped at
4{,}096 values to keep a meaningful bottleneck, did not offer the larger latent.
The comparisons in Sections~\ref{sec:task2} and~\ref{sec:task3} are read with this
in mind. The three specialists together have 11.9\,M parameters, against 9.31\,M
for the UDAE.

\begin{figure*}[t]
\centering
\resultfig{task2_classifier_optuna_intermediate_values.png}{0.32\textwidth}{From
\path{outputs/task2/optuna/task2_classifier/intermediate_values.png}.}\hfill
\resultfig{task2_specialists_optuna_optimization_history.png}{0.32\textwidth}{From
\path{outputs/task2/optuna/task2_specialists/optimization_history.png}.}\hfill
\resultfig{task2_specialists_optuna_param_importances.png}{0.32\textwidth}{From
\path{outputs/task2/optuna/task2_specialists/param_importances.png}.}
\caption{Optuna studies of Task~2. Left: validation macro-F1 per epoch of every
classifier trial; all trials approach 1.0 within a few epochs, and the median
pruner stops most of them at the fourth epoch. Centre and right: the specialist
study's optimisation history and parameter importances.}
\label{fig:task2optuna}
\end{figure*}

\subsubsection{Final training}

The classifier was retrained for 60 epochs (early-stopping patience 15) and each
specialist for 40 epochs (patience 10), independently and with its own seed,
from the committed configurations. These runs were first made on Colab and are
being repeated on Kaggle, with identical configurations, after the Colab GPU
quota ran out (Section~\ref{sec:setup}); the results below come from the Kaggle
run. On Kaggle the classifier early-stopped with its best validation macro-F1,
0.9946, at epoch 20, and the three specialists---now with one limited
$32\times32$ skip each (\texttt{configs/task2\_specialists\_skip.yaml}), mirroring
the deployed Task~1 model---reached their best validation scores at epochs 39
(salt-and-pepper), 39 (blur) and 35 (occlusion) of 40.

\subsection{Results}

\begin{table}[t]
\caption{Task 2 classifier on the test manifest: per-class precision, recall and
F1 (3{,}669 clean entries and 11{,}007 per corruption), macro averages and overall
accuracy.}
\label{tab:task2clf}
\centering
\resulttable{task2_classifier_test_per_class.tex}{%
\begin{tabular}{@{}lcccc@{}}
\toprule
Class & Precision & Recall & F1 & $n$ \\
\midrule
clean & \tbd & \tbd & \tbd & 3669 \\
salt & \tbd & \tbd & \tbd & 11007 \\
blur & \tbd & \tbd & \tbd & 11007 \\
occlusion & \tbd & \tbd & \tbd & 11007 \\
macro avg & \tbd & \tbd & \tbd & 36690 \\
accuracy & & & \tbd & 36690 \\
\bottomrule
\end{tabular}%
\missingtable{task2_classifier_test_per_class.tex}{outputs/task2/tables/classifier_test_per_class.tex}}
\end{table}

\begin{figure}[t]
\centering
\resultfig{task2_classifier_confusion_test.png}{0.85\columnwidth}{From
\path{outputs/task2/figures/classifier_confusion_test.png} (Kaggle run):
row-normalised $4\times4$ confusion matrix, rows = true class (clean, salt,
blur, occlusion), counts in brackets.}
\caption{Normalised confusion matrix of the corruption classifier on the test
manifest. Rows are true classes and sum to one, so each row is a recall profile;
this is the same normalisation used for the Task 3 routing matrix, which makes
the two directly comparable.}
\label{fig:task2confusion}
\end{figure}

The classifier makes 74 errors on the 36{,}690 test entries: accuracy
\textbf{0.998}, macro-precision 0.997, macro-recall 0.997 and macro-F1
\textbf{0.997} (Table~\ref{tab:task2clf}, Fig.~\ref{fig:task2confusion}).
Salt-and-pepper is recognised perfectly at every level, and so are the medium and
high levels of blur and occlusion. All errors involve the clean class and the
\emph{mildest} corruptions: 26 clean photographs are called blurred and 15
occluded (clean recall 0.989), 5 low-blur entries (0.14\% of that level) and 28
low-occlusion entries (0.76\%) are called clean.

\begin{figure}[t]
\centering
\resultfig{task2_classifier_blur_detection.png}{\columnwidth}{From
\path{outputs/task2/figures/classifier_blur_detection.png} (Kaggle run):
P(blur) and P(clean) against the applied blur strength $s(k,\sigma)$ with 95\%
Wilson intervals; strength 0 = clean photographs.}
\caption{Clean-versus-blur discrimination as a function of the true blur
strength. The analysis uses the applied-kernel strength rather than $\sigma$, for
the reason given in Section~\ref{sec:data}.}
\label{fig:task2blur}
\end{figure}

Both expected confusion mechanisms are visible, and both are measurable.
\emph{Soft photographs look blurred:} splitting clean photographs into quartiles
of natural sharpness (variance of the Laplacian~\cite{pech2000diatom}), the false
``blur'' rate is 2.5\% in the softest quartile, 0.3\% in the second and zero in
the two sharpest (Fig.~\ref{fig:task2blur}); the clean photographs called blurred
have a median sharpness of 0.0052 against 0.0162 for those classified correctly.
Contrary to our expectation, the classifier is not fooled by the mildest
\emph{real} blur: the low level, strength 0.65~px, is detected in 99.9\% of
cases. \emph{Dark content looks occluded:} clean photographs with at least 30\%
near-black pixels are called occluded in 5.9\% of cases, against none below 5\%;
conversely, a rectangle whose area was already at least 60\% black is detected in
only 87\% of cases, against 99.9\% when it covers bright content. Black letterbox
bands are the typical clean-to-occlusion error in Fig.~\ref{fig:task2failures}.
Both mechanisms are properties of the data, not of the model size: the evidence
for the correct class is genuinely weak.

\begin{table*}[t]
\caption{Hard-routed restoration on the test manifest: oracle routing (true label)
against predicted routing (classifier) per corruption type and severity, with
routing accuracy, PSNR and SSIM in both modes, the number of misrouted entries and
the number of harmful misroutings (predicted routing at least 1~dB worse than
oracle routing).}
\label{tab:task2routing}
\centering
\footnotesize
\resulttable{task2_routing_oracle_vs_predicted.tex}{%
\begin{tabular}{@{}llcccccccc@{}}
\toprule
Type & Level & $n$ & Route acc. & PSNR or. & PSNR pr. & SSIM or. & SSIM pr. & Misr. & Harmful \\
\midrule
clean & none & 3669 & \tbd & exact & \tbd & 1.0000 & \tbd & \tbd & \tbd \\
salt & low & 3669 & \tbd & \tbd & \tbd & \tbd & \tbd & \tbd & \tbd \\
\bottomrule
\end{tabular}%
\missingtable{task2_routing_oracle_vs_predicted.tex}{outputs/task2/tables/routing_oracle_vs_predicted.tex}}
\end{table*}

Under oracle routing the specialists reach 33.07~dB / 0.958 on salt-and-pepper,
30.35~dB / 0.909 on blur and 23.95~dB / 0.831 on occlusion, 29.12~dB / 0.899
pooled over the corrupted entries, against 27.93~dB / 0.879 for the deployed
universal model (Table~\ref{tab:crosstask}); clean inputs are returned exactly.
Predicted routing is practically identical: routing accuracy on the corrupted
entries is 0.999, the pooled PSNR differs by $-0.002$~dB, and 3{,}628 of the 3{,}669
clean inputs (98.9\%) still pass through the identity bypass exactly. With a
near-perfect classifier, the operational system loses almost nothing to its
router.

\textbf{Clean inputs and the PSNR cap.} Under oracle routing every clean input
passes through the identity bypass, so the output equals the target exactly:
zero MSE, SSIM of one, and a PSNR that is infinite and stored at the
implementation's 100~dB cap. These entries are reported as ``exact ($n/n$)''
rather than averaged as 100~dB, and the pooled restoration figures use the
\emph{corrupted} rows only. Under predicted routing the share of clean inputs
that remain exact equals the classifier's clean recall, and misrouted clean
photographs receive a finite PSNR that is reported separately.

\subsection{Failure analysis}

\begin{figure}[t]
\centering
\resultfig{task2_routing_failures.png}{\columnwidth}{From
\path{outputs/task2/figures/routing_failures.png} (Kaggle run): worst harmful
misroutings per confusion pair, showing target, input, oracle output, predicted
output and the four classifier probabilities.}
\caption{Cases where a classifier error causes a restoration failure.}
\label{fig:task2failures}
\end{figure}

\begin{table}[t]
\caption{Every classifier confusion on the test manifest: number of entries, share
of the true type, mean PSNR and SSIM change of predicted relative to oracle
routing, and how many misroutings were harmful, neutral or beneficial
($\pm1$~dB).}
\label{tab:task2misrouting}
\centering
\resulttable{task2_routing_misrouting.tex}{%
\begin{tabular}{@{}llccc@{}}
\toprule
True & Routed to & $n$ & $\Delta$PSNR & Harmful / neutral / benef. \\
\midrule
clean & blur & \tbd & \tbd & \tbd \\
\bottomrule
\end{tabular}%
\missingtable{task2_routing_misrouting.tex}{outputs/task2/tables/routing_misrouting.tex}}
\end{table}

Table~\ref{tab:task2misrouting} lists the four confusions. \emph{Clean routed to
an expert} (26 to blur, 15 to occlusion) is always harmful by construction, since
the oracle output is exact; the blur specialist still returns these images at
33.0~dB on average, the occlusion specialist at 24.6~dB because it ``fills'' dark
bands that were content (Fig.~\ref{fig:task2failures}). \emph{Low occlusion routed
to the identity} (28 entries) is the most damaging confusion: the rectangle is
left in place and PSNR falls by 3.05~dB on average (23 harmful, 3 neutral, 2
beneficial). \emph{Low blur routed to the identity} (5 entries) is the opposite
case: 4 of the 5 are \emph{beneficial} and the mean change is $+1.47$~dB, because
for a blur this mild the unmodified input is closer to the clean image than the
specialist's reconstruction. Misrouting therefore matters only at the low
severities, and not always in the harmful direction.

\section{Task 3: Soft Mixture-of-Experts Restoration}
\label{sec:task3}

\subsection{Methodology}

\subsubsection{Model}

The soft mixture replaces the \texttt{argmax} of Task~2 with a temperature-scaled
softmax over the same four branches:
\begin{align}
  w &= \mathrm{softmax}\!\left(\frac{G(\tilde{x})}{\tau}\right), \label{eq:gate}\\
  \hat{x} &= w_0\,\tilde{x} + w_1 A_\text{salt}(\tilde{x}) + w_2 A_\text{blur}(\tilde{x}) + w_3 A_\text{occ}(\tilde{x}).
  \label{eq:moe}
\end{align}
The gate $G$ is the Task~2 classifier and the three experts are the Task~2
specialists, loaded from their checkpoints and rebuilt from their own stored
configurations, so whatever architecture the Task~2 search selected is used
unchanged. Branch~0 is a \emph{true identity}, not a learned clean expert: a
clean input needs only the gate to place its weight on branch~0 to be reproduced
perfectly, so there is no reconstruction floor for clean images.

The temperature $\tau$ is a registered buffer rather than a parameter or a Python
float, which means it is saved with the state dictionary, is never touched by the
optimiser, and is traced into the exported ONNX graph as a constant, so the
deployed model provably uses the temperature it was trained with. Temperature
scaling of a softmax follows Hinton et al.~\cite{hinton2015distilling}: a higher
$T$ produces a softer distribution. Here a small $\tau$ approaches the hard
routing of Task~2 and a large $\tau$ approaches uniform blending, so $\tau$ is
the single knob that interpolates between Tasks~2 and~3---which is why it is
searched. (A stochastic hard route with a differentiable relaxation, as in the
Gumbel-softmax~\cite{jang2017categorical}, was the alternative not taken: the
brief asks for a continuous weight on every branch, and a dense blend needs no
relaxation.) The gate logits and the softmax are computed in float32 even under
mixed precision, since the routing decision multiplies \emph{every} branch
output.

\textbf{A caveat that should be stated plainly.} Equation~(\ref{eq:moe}) trained
on the error of the blend is exactly the \emph{cooperative} error function of
Jacobs et al.~\cite{jacobs1991adaptive}, the formulation they argue ``does not
encourage localization'' because each expert need only cancel the residual left
by the others. Two features of this system counteract that. The experts are
pre-trained to specialise rather than learned from scratch, and the
cross-entropy term below supervises the routing directly with the known runtime
corruption label. Without both, the blend objective would be expected to drift
towards using many experts for every input.

\subsubsection{Loss}

\begin{equation}
  \mathcal{L}_\text{MoE} = \lambda_1 \mathcal{L}_{L_1} + \lambda_s (1 - \mathrm{SSIM}) + \lambda_c \mathcal{L}_\text{CE} + \lambda_b \mathcal{L}_\text{balance}.
  \label{eq:moeloss}
\end{equation}
The cross-entropy is applied to the \emph{raw} gate logits $G(\tilde{x})$ rather
than to $G(\tilde{x})/\tau$. With the scaled logits, the same parameter $\tau$
would control both the sharpness of the routing and the magnitude of the
classification gradient---a small $\tau$ multiplies the cross-entropy gradient by
$1/\tau$---so the searches over $\tau$ and $\lambda_c$ would be confounded.
Keeping the cross-entropy on the raw logits makes $\tau$ purely a routing knob,
$\lambda_c$ purely a weight, and preserves the Task~2 objective so that the gate
remains a calibrated corruption classifier whose weights are interpretable in the
gating analysis and in the application display.

\subsubsection{The balance regulariser}

The regulariser is the brief's,
\begin{equation}
  \mathcal{L}_\text{balance} = \sum_{k=0}^{3} \Bigl( \bar{w}_k - \tfrac{1}{4} \Bigr)^2,
  \label{eq:balance}
\end{equation}
with $\bar{w}_k$ the mean weight of branch $k$ over the batch. Three arguments
support it, and the alternatives were examined in detail before it was kept.

\emph{First, it is the appropriate simplification for a dense gate.} Shazeer et
al.\ use two terms~\cite{shazeer2017outrageously}: an importance loss
$\mathcal{L}_\text{imp} = w_\text{imp} \cdot \mathrm{CV}(\mathrm{Importance}(X))^2$,
the squared coefficient of variation of the batchwise sum of gate values, and a
load loss $\mathcal{L}_\text{load} = w_\text{load} \cdot \mathrm{CV}(\mathrm{Load}(X))^2$
over a smooth estimator of the \emph{number of examples} routed to each expert.
The Switch Transformer merges them into
\begin{equation}
  \mathcal{L}_\text{switch} = \alpha\,N \sum_{i=1}^{N} f_i P_i,
  \label{eq:switch}
\end{equation}
where $f_i$ is the fraction of tokens dispatched to expert $i$ and $P_i$ the mean
router probability, with $\alpha = 10^{-2}$ chosen after a sweep over five orders
of magnitude~\cite{fedus2022switch,lepikhin2021gshard}. Both the load term and
$f_i$ exist to count \emph{discrete dispatches} in a top-$k$ router. The gate
here is dense: all four branches always run, nothing is dropped, and there is no
load to count, only importance. With soft routing $f$ and $P$ coincide, and since
$\sum_k (P_k - 1/N)^2 = \sum_k P_k^2 - 1/N$, equation~(\ref{eq:balance}) is the
Switch loss up to an affine transformation. The brief's term is therefore not a
weaker substitute for the published losses; for a dense gate it is the same
objective.

\emph{Second, it does not fight the supervised objective.} Because training
batches contain exactly $B/4$ images of each condition, perfect routing---one-hot
on the true class---gives every branch a mean weight of exactly $1/4$, so
$\mathcal{L}_\text{balance} = 0$ at the optimum of the cross-entropy. The two
terms are compatible by construction, a property that follows from the balanced
sampler and is unavailable in the MoE literature, which has no labels. (With an
unbalanced batch the term would actively punish correct routing, which a unit
test asserts.) Total collapse onto a single branch costs
$(1 - 1/4)^2 + 3 \cdot (1/4)^2 = 0.75$.

\emph{Third, the exact form matters less than its presence.} Shazeer et al.'s own
ablation, reproduced in Table~\ref{tab:shazeer}, shows that any combination
containing at least one balance term gives nearly identical model quality over
two orders of magnitude of weight, whereas omitting them entirely is far worse.

\begin{table}[t]
\caption{Effect of the balance losses in Shazeer et al.~\cite{shazeer2017outrageously},
Table 6 (MoE-256 model, 10 epochs). Any balance term suffices; none is much
worse. $\mathrm{CV}$ denotes the coefficient of variation.}
\label{tab:shazeer}
\centering
\begin{tabular}{@{}cccccc@{}}
\toprule
$w_\text{imp}$ & $w_\text{load}$ & Perplexity & $\mathrm{CV}(\text{Imp.})$ & $\mathrm{CV}(\text{Load})$ & $\frac{\max}{\text{mean}}$ \\
\midrule
0.0 & 0.0 & 39.8 & 3.04 & 3.01 & 17.80 \\
0.2 & 0.0 & 35.6 & 0.06 & 0.17 & 1.47 \\
0.0 & 0.2 & 35.7 & 0.22 & 0.04 & 1.15 \\
0.1 & 0.1 & 35.6 & 0.06 & 0.05 & 1.14 \\
0.01 & 0.01 & 35.7 & 0.48 & 0.11 & 1.37 \\
1.0 & 1.0 & 35.7 & 0.03 & 0.02 & 1.07 \\
\bottomrule
\end{tabular}
\end{table}

Entropy regularisers were considered and rejected as a \emph{loss}. Maximising
routing entropy would push against the cross-entropy by rewarding weight spread
onto wrong branches; minimising it would push towards hard routing, which is
Task~2. (The information-maximisation combination of a high marginal entropy and
a low conditional entropy~\cite{krause2010discriminative,grandvalet2004semi}
remains the principled upgrade if a stronger regulariser is ever needed, and
confidence penalties~\cite{pereyra2017regularizing} the remedy if the gate
saturates.) Entropy is used here as a \emph{diagnostic} instead, reported per
class. The one blind spot of~(\ref{eq:balance}) is that it sees only batch-mean
weights and is therefore insensitive to a gate that routes each class
confidently to the \emph{wrong} branch; the cross-entropy term and the routing
matrix below exist for that.

\subsubsection{Routing collapse}

Collapse is defined once and used identically by the Optuna pruner, the per-epoch
log and the final evaluation. From the routing matrix on the balanced validation
manifest, $R[c,k]$ being the mean weight of branch $k$ over entries of true class
$c$ (rows sum to one), define $\mathrm{usage}_k = \mathrm{mean}_c R[c,k]$, which
is $1/4$ under balanced perfect routing, and
$\mathrm{off}_k = \mathrm{mean}_{c \neq k} R[c,k]$. A configuration has collapsed
if some branch is \emph{dead} ($\mathrm{usage}_k < 0.05$, a fifth of its fair
share; since $\mathrm{usage}_k \geq R[k,k]/4$ such a branch receives less than
$0.2$ of the weight even on the inputs it exists for) or \emph{dominating}
($\mathrm{off}_k > 0.5$, taking more than half the weight on average on the other
three classes). The criterion deliberately does not flag a merely uncertain
model: uniform routing is informative, not useless, and the tests assert that
uniform weights are \emph{not} flagged while total collapse is flagged as both
dead and dominating.

\subsubsection{Staged training}

Training proceeds in the two stages the brief prescribes. In the \textbf{warm-up}
the experts are frozen in two senses at once: their parameters have
\texttt{requires\_grad = False} and the warm-up optimiser is built over the gate
only, and they are kept in evaluation mode with the \texttt{train()} method
overridden so that a later call cannot silently undo it. The second part is not
redundant---without it, BatchNorm would keep updating its running statistics on
every forward pass and dropout would stay active, so the ``frozen'' experts would
drift and the warm-up would no longer be a clean measurement of the gate. In the
\textbf{joint stage} all parameters train, with the gate at learning rate
$\eta$ and the experts at $0.1\eta$, under cosine decay.

The schedule's shape has direct empirical support. Kumar et
al.~\cite{kumar2022finetuning} show that full fine-tuning from a randomly
initialised head distorts good pretrained features, because the lower layers move
while the head is still wrong, and that linear probing followed by full
fine-tuning combines the benefits of both. The warm-up is exactly the probe
stage, with the gate as the head: it prevents a miscalibrated gate from pushing
destructive gradients into three specialists that are already good. The lower
expert learning rate in the joint stage is discriminative
fine-tuning~\cite{howard2018universal}, and is additionally motivated by the fact
that the gradient reaching an expert is scaled by its routing weight, so a
rarely selected expert receives little signal while a dominant one receives a
lot. The experts' BatchNorm running statistics are kept fixed during joint
training, since an expert's input distribution changes once it sees all four
conditions and re-estimating would move its normalisation away from the regime it
was trained in~\cite{he2019rethinking}.

The validation objective caps per-image PSNR at 40~dB \emph{inside the objective
only}, because a clean input routed to the identity reaches 60--100~dB and would
otherwise dominate model selection over actual restoration quality. The uncapped
score is logged, and reported test metrics are never capped this way.

\subsubsection{Optuna search design}

The study \texttt{task3\_moe} searches the five parameters the brief names
(Table~\ref{tab:task3search}), with trial~0 enqueued at the brief's starting
values $\lambda_1 = 0.8$, $\lambda_s = 0.2$, $\lambda_c = 0.1$,
$\lambda_b = 0.01$. Pruning has two paths, both recorded in the trial
attributes: the median pruner for weak configurations, and the collapse criterion
above---a task-specific stopping rule that no generic pruner can express, and
which the brief explicitly permits. A diverged loss is converted into a prune with
its reason rather than a crashed trial, so one bad learning rate cannot end a
session.

\textbf{Budget actually used.} The study was planned as 20 trials of one warm-up
and four joint epochs, followed by a final run of 2 warm-up and 20 joint epochs.
Because the Colab GPU quota ran out on the deadline day, Task~3 was trained on
Kaggle with a reduced budget: \textbf{6 trials of one warm-up and three joint
epochs} (25-minute cap), and a \textbf{final run of 2 warm-up and 10 joint
epochs}. Six trials are the brief's reference configuration plus five random
proposals---the TPE sampler only starts modelling after ten random start-up
trials---so this search is a small random search around the reference point
rather than a model-based optimisation, a limitation stated in
Section~\ref{sec:limitations}.

\begin{table}[t]
\caption{Optuna search space for Task 3 (study \texttt{task3\_moe}); trial~0 is the
brief's starting configuration.}
\label{tab:task3search}
\centering
\footnotesize
\setlength{\tabcolsep}{3pt}
\begin{tabular}{@{}llp{1.6cm}c@{}}
\toprule
Parameter & Range & Rationale & Selected \\
\midrule
Joint learning rate & $10^{-5}$--$3\times10^{-4}$ & log; capped at the warm-up rate & $3.57\times10^{-5}$ \\
Temperature $\tau$ & $0.3$--$3.0$ & log; near-hard to near-uniform & 2.68 \\
$\lambda_c$ & $0.01$--$1.0$ & log; a decade either side of 0.1 & 0.291 \\
$\lambda_b$ & $10^{-3}$--$10^{-1}$ & log; contains Switch's $10^{-2}$ & 0.0158 \\
$\alpha$ ($=\lambda_1 = 1-\lambda_s$) & $0.5$--$0.95$ & constant reconstruction scale & 0.570 \\
\bottomrule
\end{tabular}
\\[2pt]
Six trials ran in 19.3 minutes (five completed, one pruned by the median rule).
The best was trial~1 with a validation score of 0.8671, against 0.8667 for the
brief's reference configuration (trial~0); all six trials lie within 0.8661--0.8671.
The mixture starts from trained components, so over four epochs its score is almost
insensitive to these hyperparameters---including $\tau$, which varied from 0.34 to
2.80 across trials.
\end{table}

\begin{figure*}[t]
\centering
\resultfig{task3_optuna_optimization_history.png}{0.42\textwidth}{From
\path{outputs/task3/optuna/task3_moe/optimization_history.png}.}\hfil
\resultfig{task3_optuna_intermediate_values.png}{0.42\textwidth}{From
\path{outputs/task3/optuna/task3_moe/intermediate_values.png}.}
\caption{Optuna study \texttt{task3\_moe} (six trials on Kaggle): optimisation
history (left) and validation score per epoch of every trial, including pruned
ones (right).}
\label{fig:task3optuna}
\end{figure*}

\subsection{Results}

The selected configuration was fine-tuned for 2 warm-up and 10 joint epochs,
starting from the new Task~2 checkpoints (classifier epoch 20; specialists epochs
39, 39 and 35); the best validation score, 0.8673, was reached at epoch~8, and the
routing-collapse check never fired. On the test manifest the mixture reaches
29.45~dB / 0.902 on the corrupted entries (salt-and-pepper 33.62~dB / 0.961, blur
30.72~dB / 0.912, occlusion 24.02~dB / 0.834) and 60.77~dB / 0.999 on clean inputs
(Table~\ref{tab:task3results}).

\begin{table}[t]
\caption{Soft mixture-of-experts restoration on the test manifest, per input type
and severity. PSNR* excludes outputs identical to the target (PSNR capped at
100~dB); ``Exact'' counts them.}
\label{tab:task3results}
\centering
\resulttable{task3_test_by_type_level.tex}{%
\begin{tabular}{@{}llccccc@{}}
\toprule
Type & Level & $n$ & PSNR & PSNR* & SSIM & Exact \\
\midrule
clean & none & 3669 & \tbd & \tbd & \tbd & \tbd \\
salt & low & 3669 & \tbd & \tbd & \tbd & \tbd \\
\bottomrule
\end{tabular}%
\missingtable{task3_test_by_type_level.tex}{outputs/task3/tables/test_by_type_level.tex}}
\end{table}

\begin{table*}[t]
\caption{Cross-task comparison on identical test entries: the universal model
(Task~1), hard routing with oracle and predicted routes (Task~2) and the soft
mixture (Task~3). PSNR* excludes outputs identical to the target, so the clean row
of hard routing is often empty; the clean row is best compared by SSIM.}
\label{tab:crosstask}
\centering
\footnotesize
\resulttable{crosstask_by_type.tex}{%
\begin{tabular}{@{}lcccccccc@{}}
\toprule
\multirow{2}{*}{Input} & \multicolumn{2}{c}{Task 1 universal} & \multicolumn{2}{c}{Task 2 oracle} & \multicolumn{2}{c}{Task 2 predicted} & \multicolumn{2}{c}{Task 3 soft MoE} \\
\cmidrule(lr){2-3}\cmidrule(lr){4-5}\cmidrule(lr){6-7}\cmidrule(lr){8-9}
 & PSNR* & SSIM & PSNR* & SSIM & PSNR* & SSIM & PSNR* & SSIM \\
\midrule
clean & \tbd & \tbd & \tbd & \tbd & \tbd & \tbd & \tbd & \tbd \\
salt & \tbd & \tbd & \tbd & \tbd & \tbd & \tbd & \tbd & \tbd \\
blur & \tbd & \tbd & \tbd & \tbd & \tbd & \tbd & \tbd & \tbd \\
occlusion & \tbd & \tbd & \tbd & \tbd & \tbd & \tbd & \tbd & \tbd \\
\bottomrule
\end{tabular}%
\missingtable{task3_comparison_by_type.tex}{outputs/task3/tables/comparison_by_type.tex}}
\end{table*}

\begin{figure*}[t]
\centering
\resultfig{task3_comparison_psnr.png}{0.49\textwidth}{From
\path{outputs/task3/figures/comparison_psnr.png}.}\hfill
\resultfig{task3_comparison_ssim.png}{0.49\textwidth}{From
\path{outputs/task3/figures/comparison_ssim.png}.}
\caption{Cross-task comparison per input type and severity: PSNR* (left) and SSIM
(right) of the universal model, hard routing (oracle and predicted) and the soft
mixture on identical test entries.}
\label{fig:crosstask}
\end{figure*}

\begin{table}[t]
\caption{Paired per-entry comparison on the corrupted test entries (built by
\texttt{report/crosstask\_table.py} from the four per-entry record files): mean PSNR
gain and share of entries on which the first system is better.}
\label{tab:paired}
\centering
\footnotesize
\resulttable{crosstask_paired.tex}{\missingtable{crosstask_paired.tex}{report/crosstask_table.py}}
\end{table}

\textbf{Cross-task comparison.} Tables~\ref{tab:crosstask} and~\ref{tab:paired}
compare the deployed universal model (\texttt{udae\_skip32}), hard routing and the
soft mixture on identical entries. The ranking is the same for every corruption:
soft mixture $>$ hard routing $>$ universal model. Over the 33{,}021 corrupted
entries the mixture reaches 29.45~dB / 0.902, hard routing 29.12~dB / 0.899 in
both modes and the universal model 27.93~dB / 0.879. Paired, the mixture beats the
universal model on 90.7\% of the corrupted entries (+1.52~dB on average; on all
11{,}007 salt-and-pepper entries) and hard predicted routing on 83.3\% (+0.33~dB;
99.8\% of salt-and-pepper, 93.8\% of blur but only 56.3\% of occlusion entries,
+0.08~dB). Hard routing beats the universal model on 84.8\% of entries
(+1.19~dB). The gains are largest for salt-and-pepper and smallest for
occlusion, where every system is limited by the missing content. Clean inputs
separate the designs most clearly: the universal model returns them at 32.21~dB,
hard routing exactly in 3{,}628 of 3{,}669 cases, and the mixture at 60.77~dB with
SSIM 0.9992. Two caveats apply: the specialists' latent is twice the universal
model's (Section~\ref{sec:task2}), and the specialists and the deployed universal
model both use one $32\times32$ skip, so the comparison isolates routing from
the skip but not from latent size. The soft mixture's advantage over hard
routing does not come from blending (the routing below is nearly one-hot) but
from the joint fine-tuning of gate and experts on all four conditions.

\subsubsection{Gating analysis}

\begin{figure}[t]
\centering
\resultfig{task3_routing_heatmap_type_level.png}{\columnwidth}{From
\path{outputs/task3/figures/routing_heatmap_type_level.png}: mean weight of the
four branches (columns) for every true type and severity (rows).}
\caption{Average expert weights for every true corruption type and severity
level. Rows sum to one.}
\label{fig:task3heatmap}
\end{figure}

\begin{figure}[t]
\centering
\resultfig{task3_weight_distributions.png}{\columnwidth}{From
\path{outputs/task3/figures/weight_distributions.png}: distribution of each branch
weight within each true class.}
\caption{Routing-weight distributions per true class.}
\label{fig:task3weights}
\end{figure}

\begin{table}[t]
\caption{Expert activity on the test set: class-balanced usage of each branch
(1/4 under balanced perfect routing), mean weight on its own class and on the
other classes, the largest weight on any other class, and how often it dominates
other classes.}
\label{tab:task3activity}
\centering
\resulttable{task3_expert_activity.tex}{%
\begin{tabular}{@{}lccccc@{}}
\toprule
Branch & Usage & Own & Others & Max other & on \\
\midrule
identity & \tbd & \tbd & \tbd & \tbd & \tbd \\
\bottomrule
\end{tabular}%
\missingtable{task3_expert_activity.tex}{outputs/task3/tables/expert_activity.tex}}
\end{table}

\begin{figure*}[t]
\centering
\resultfig[0.24\textheight]{task3_gating_dominant.png}{0.49\textwidth}{From
\path{outputs/task3/figures/gating_dominant.png}.}\hfill
\resultfig[0.24\textheight]{task3_gating_spread.png}{0.49\textwidth}{From
\path{outputs/task3/figures/gating_spread.png}.}
\caption{Inputs on which one expert dominates (left, lowest routing entropy per
true class) and inputs on which the weights are spread across several branches
(right, highest entropy), each with all four branch outputs and their weights,
the mixed output and the target. Examples are selected by entropy rather than by
hand, so they are reproducible.}
\label{fig:task3gating}
\end{figure*}

\textbf{Routing is sharp, not spread.} Despite the selected temperature
$\tau = 2.68$, near the soft end of the range, the gate routes almost one-hot: the
largest weight is on the correct branch for 99.8\% of test entries, the mean
routing entropy normalised by $\log 4$ is 0.074, and the diagonal of the routing
matrix is 0.936 (clean to identity), 0.996, 0.947 and 0.989
(Fig.~\ref{fig:task3heatmap}, Table~\ref{tab:task3activity}). No branch is dead or
dominating: usage is 0.250, 0.251, 0.247 and 0.252 against the ideal 0.25, and the
largest mean weight any branch takes on other classes is 0.022, so the collapse
criterion is far from triggering. The likely explanation is that the cross-entropy
term acts on the raw gate logits, initialised from a classifier trained to
confidence, so the logits are large and dividing them by 2.68 still leaves
near-one-hot weights: $\tau$ and the logit scale trade off, which is also why the
score barely depended on $\tau$ in the search. Softness appears exactly where the
evidence is weakest: the entropy is highest for clean inputs (0.170) and blur
(0.148, 0.184 at the low level), where the gate hedges between the identity and the
blur expert---a clean image receives 0.041 blur weight on average and a low blur
0.071 identity weight---while salt-and-pepper (0.016) and occlusion (0.026) are
routed almost deterministically; mild corruptions are routed more softly than
strong ones. The high-entropy examples in Fig.~\ref{fig:task3gating} are of this
kind: dark or soft images, and small rectangles on dark regions, where identity
and the relevant expert share the weight.

\subsection{Failure analysis}

\begin{figure}[t]
\centering
\resultfig[0.24\textheight]{task3_failures_misrouted.png}{\columnwidth}{From
\path{outputs/task3/figures/failures_misrouted.png}: the six lowest-SSIM entries
whose largest weight is not on the true branch.}
\caption{Task 3 failure cases: mis-routed entries with the lowest SSIM.}
\label{fig:task3failures}
\end{figure}

\begin{table}[t]
\caption{Beyond the brief: two corruptions per image (medium severity). Mean
mixture weights per combination and the share of entries whose two largest
weights are exactly the two corruptions present.}
\label{tab:task3mixed}
\centering
\resulttable{task3_mixed_routing.tex}{%
\begin{tabular}{@{}lcccccc@{}}
\toprule
Input & $n$ & Identity & Salt & Blur & Occl. & Top-2 = parts \\
\midrule
salt+blur & \tbd & \tbd & \tbd & \tbd & \tbd & \tbd \\
\bottomrule
\end{tabular}%
\missingtable{task3_mixed_routing.tex}{outputs/task3/mixed/mixed_routing.tex}}
\end{table}

The worst entries whose largest weight is on a wrong branch
(Fig.~\ref{fig:task3failures}) are the cases Task~2 also missed---small rectangles on
dark regions and soft clean photographs---so soft routing recovers misroutings
only partly: blending the identity with the right expert leaves part of a rectangle
or part of the blur. On average, however, the joint fine-tuning pays off: the
mixture beats hard predicted routing on 83.3\% of the corrupted entries
(Table~\ref{tab:paired}).

\textbf{Beyond the brief: two corruptions per image.} As an extra experiment,
clearly outside the brief's corruption definitions, 900 test images were given two
corruptions at medium severity (Table~\ref{tab:task3mixed}). One might expect the
gate to split its weight between the two corresponding experts; it does not. It
picks a single expert with a fixed precedence---salt-and-pepper over blur and over
occlusion, occlusion over blur---with weight 0.998 or more, and the two largest
weights are exactly the two corruptions present only for salt plus occlusion, where
the second weight is negligible. Consequently soft and argmax routing of the same
mixture are indistinguishable (20.80~dB both, against 13.37~dB for the input), and
the mixture is only 0.12~dB better than Task~2's hard routing (20.69~dB). For salt
plus occlusion every system stays near 13.4~dB, since the salt-and-pepper expert
removes the impulses but leaves the rectangle. Restoring compound degradations
would need experts trained on them or genuinely sequential routing.

\section{Task 4: Style-Conditioned Face-to-Sketch Generation}
\label{sec:task4}

\subsection{Methodology}

\subsubsection{Generator}

The generator follows the pix2pix \texttt{unet\_128} layout~\cite{isola2017image}:
seven $4\times4$ stride-2 convolutions down to $1\times1$ with LeakyReLU$(0.2)$,
then six transposed convolutions with skip connections and a $\tanh$ output,
with dropout after the three innermost decoder layers, following the
convolution-normalisation-activation modules that pix2pix inherits from
DCGAN~\cite{radford2016unsupervised} together with its $\mathcal{N}(0, 0.02)$
weight initialisation. It is worth noting the
contrast with Section~\ref{sec:task1}: the U-Net skips that are \emph{forbidden}
there, because they would defeat the bottleneck, are \emph{required} here,
because the goal is a translation that preserves the input's spatial layout
rather than a compressed representation. Architectural choices follow the
objective.

Instance normalisation is used rather than batch normalisation. pix2pix's batch
normalisation with batch size one and test-batch statistics is instance
normalisation in all but name~\cite{isola2017image,ulyanov2016instance}; making
it explicit means training and evaluation behave identically---so ONNX inference
matches training---and makes the network independent of the batch size, which
Optuna tunes.

\textbf{Style conditioning} uses a single learned embedding of the three styles.
At each of the six decoder layers, a conditional instance normalisation computes
$\mathrm{IN}(h) \cdot (1 + \gamma_l(e)) + \beta_l(e)$, where $\gamma_l$ and
$\beta_l$ are linear projections of the shared embedding $e$. This is FiLM in
the form of conditional instance
normalisation~\cite{perez2018film,dumoulin2017learned,dumoulin2018featurewise};
projecting one shared embedding to every layer's gains and biases follows
BigGAN's shared embedding~\cite{brock2019large}, and modulating the decoder is
what StarGAN~v2 does for multi-domain translation~\cite{choi2020starganv2,huang2017arbitrary}.
The encoder stays style-agnostic, since it only has to read the photograph.

The alternatives were not merely argued but measured in a controlled experiment:
one synthetic photograph with three target sketches differing only in style, the
same U-Net, two seeds, comparing the $L_1$ error under the true style against
the error under a shifted style. A style-blind model can only reach the median of
the three targets, around $0.054$. Input concatenation scored $0.0558$ with a
gap of \emph{exactly zero}---it learned nothing from the label, as the theory
predicts: a spatially constant input map is only a conditional bias on the first
layer, and the instance normalisation of the next layer subtracts per-channel
offsets, so the signal is normalised away~\cite{dumoulin2018featurewise}.
Injecting at the $1\times1$ bottleneck gave a gap of $0.0026$, limited because
the U-Net skips let the decoder bypass the bottleneck entirely. Decoder FiLM gave
$0.0452$ against $0.0555$, a gap of $0.0104$---four times the bottleneck
variant---and also modulating the encoder added nothing measurable ($0.0108$).
Decoder FiLM was therefore chosen.

\subsubsection{Discriminator}

The discriminator is the pix2pix $70\times70$ PatchGAN on the concatenation of
photograph and sketch, with three downsampling layers, instance normalisation and
LeakyReLU$(0.2)$, producing a $14\times14$ map of logits for $128\times128$
inputs. The $70\times70$ receptive field was verified by gradient propagation
with the normalisations disabled. pix2pix found this patch size best among 1, 16,
70 and 286 pixels at $256\times256$ resolution: a $16\times16$ patch is enough for
sharpness but tiles, and the full-image discriminator scored considerably
worse~\cite{isola2017image}. At $128\times128$ a 70-pixel patch covers about half
a face.

\textbf{Style conditioning in the discriminator is by projection}: the output
layer is a shared $4\times4$ kernel plus a style-specific kernel generated from
the discriminator's own embedding, so each patch logit is
$\mathrm{logit}(p) = w \cdot \phi(p) + b + v(s) \cdot \phi(p)$. This is Miyato
and Koyama's projection discriminator,
$f(x,y) = y^{\mathsf{T}} V \phi(x) + \psi(\phi(x))$, applied per
patch~\cite{miyato2018cgans}; it follows from assuming that $p(y \mid x)$ is
log-linear in the features, and on ImageNet it beat both concatenation and the
AC-GAN auxiliary classifier~\cite{odena2017conditional}. Concatenation would face
the same normalisation problem as in the generator. The photograph, being a
genuinely spatial condition, is still supplied by concatenation as in pix2pix.
Implementation detail: one convolution produces a logit map per style and a
\texttt{gather} selects each row's style, so only the used embedding rows receive
gradients.

Two options are held in reserve should training prove unstable or style control
weak: spectral normalisation of the discriminator~\cite{miyato2018spectral} if it
overpowers the generator, and mismatched-condition negatives---a real sketch
presented with the wrong style and labelled fake, as in
GAN-CLS~\cite{reed2016generative}---if the source/style confound of
Section~\ref{sec:data} shows up as weak style control.

\subsubsection{Objective and training}

The generator objective is adversarial plus reconstruction,
\begin{equation}
  \mathcal{L}_G = \mathcal{L}_\text{adv} + \lambda_{L_1}\,\mathcal{L}_{L_1}\bigl(y, G(x,s)\bigr),
  \label{eq:ganloss}
\end{equation}
with binary cross-entropy with logits for the adversarial terms. pix2pix uses
$\lambda_{L_1} = 100$ and reports that $L_1$ alone gives blurry results while the
conditional GAN alone gives sharp results with artefacts, the combination
reducing them~\cite{isola2017image}; 100 is the starting point here and
$\lambda_{L_1}$ is searched. Following pix2pix, Adam~\cite{kingma2015adam} is
used with $\beta = (0.5, 0.999)$ and the discriminator objective is halved to
slow it relative to the generator; separate learning rates for the two networks
are standard~\cite{heusel2017gans}. All spatial augmentation is applied
identically to both members of a pair, as required to preserve pixel
correspondence. Unlike pix2pix, dropout is active only in training mode: the
application needs deterministic sketches and exact ONNX parity, so inference runs
in evaluation mode.

\subsubsection{Optuna search design}

The study \texttt{task4\_cgan} searches the seven quantities the brief names
(Table~\ref{tab:task4search}). The generator and discriminator share one base
channel count and one style-embedding dimension, which keeps their capacities
balanced and the search space small. Trial~0 is the pix2pix and brief reference
point. The objective is the best validation \emph{sketch score},
$0.5\,\mathrm{SSIM} + 0.5\,(1 - L_1)$ on the validation split only, which involves
neither $\lambda_{L_1}$ nor the discriminator, so trials with different loss
weights remain comparable; the sharpness diagnostic introduced below is recorded
for every trial so that the blur a large $\lambda_{L_1}$ induces stays visible.
TPE is seeded at 42 and the median pruner waits for four trials and three epochs.

The final run trains the selected configuration from scratch for \textbf{200
epochs}---the pix2pix setting for datasets of this size~\cite{isola2017image}---with
the pix2pix schedule: a constant learning rate for the first half, then linear
decay to zero. The search ran \textbf{18 trials of 12 epochs} (35-minute cap). A
first, smaller search (10 trials of 8 epochs) had selected a discriminator learning
rate 4.5$\times$ below the generator's, and its 120-epoch model drew soft, faint
strokes---a symptom of a discriminator too weak to penalise blur. The
discriminator learning-rate range was therefore narrowed to start at
$1.5\times10^{-4}$ and the task was retrained; only the second run is reported.

\begin{table}[t]
\caption{Optuna search space for Task 4 (study \texttt{task4\_cgan}, 18 trials of
12 epochs on Kaggle), the reference trial~0, and the selected configuration.}
\label{tab:task4search}
\centering
\setlength{\tabcolsep}{3.5pt}
\setlength{\tabcolsep}{2.5pt}\begin{tabular}{@{}lllcc@{}}
\toprule
Parameter & Range & Scale & Trial 0 & Selected \\
\midrule
$\eta_G$ & $5\times10^{-5}$--$5\times10^{-4}$ & log & $2\times10^{-4}$ & $4.04\times10^{-4}$ \\
$\eta_D$ & $1.5\times10^{-4}$--$5\times10^{-4}$ & log & $2\times10^{-4}$ & $2.03\times10^{-4}$ \\
Batch size & $\{4, 8, 16\}$ & cat. & 8 & 8 \\
Base channels & $\{32, 48, 64\}$ & cat. & 64 & 48 \\
Dropout & $0.0$--$0.5$ & uniform & 0.5 & 0.465 \\
Style dim. & $\{8, 16, 32\}$ & cat. & 16 & 32 \\
$\lambda_{L_1}$ & $25$--$200$ & log & 100 & 133 \\
\bottomrule
\end{tabular}
\\[2pt]
Eighteen trials ran in 31.0 minutes of trial time (10 completed, 8 pruned). The
best was trial~9, with a validation sketch score of 0.6710 against 0.6597 for the
pix2pix reference (trial~0). The four best trials all used $\lambda_{L_1}$
between 98 and 184 and a generator learning rate 1.5 to 2.6 times the
discriminator's.
\end{table}

The ranges are an octave either side of pix2pix's $2\times10^{-4}$ for both
learning rates, searched separately because their ratio sets the balance between
the two networks; batch sizes giving 224, 112 or 56 updates per epoch on the 899
training pairs; base widths giving 10.5, 23.6 or 41.9\,M generator parameters; a
dropout range that includes pix2pix's 0.5, its generator's only noise source, and
zero; and a $\lambda_{L_1}$ range centred on the brief's 100.

\begin{figure*}[t]
\centering
\resultfig{task4_optuna_optimization_history.png}{0.42\textwidth}{From
\path{outputs/task4/optuna/task4_cgan/optimization_history.png}.}\hfil
\resultfig{task4_optuna_param_importances.png}{0.42\textwidth}{From
\path{outputs/task4/optuna/task4_cgan/param_importances.png}.}
\caption{Optuna study \texttt{task4\_cgan} (ten trials on Kaggle): optimisation
history (left) and parameter importances (right).}
\label{fig:task4optuna}
\end{figure*}

\subsection{Results}

The selected configuration was trained for 200 epochs; the best validation
sketch score, 0.6885, was reached at epoch 97. The losses stayed finite under fp16
autocast throughout---the training code aborts on any non-finite loss, and the run
completed.

\begin{table}[t]
\caption{Task 4 results on the official FS2K test set (1{,}046 pairs), per style:
$L_1$ (on $[0,1]$), SSIM, PSNR (dB), edge ratio (generated edge energy divided by
the ground truth's; 1 = as sharp as the artist) and sketch score. Note the
support: style 3 has only 46 test pairs.}
\label{tab:task4}
\centering
\resulttable{task4_by_style.tex}{%
\begin{tabular}{@{}lcccccc@{}}
\toprule
Style & $n$ & $L_1$ & SSIM & PSNR & Edge ratio & Score \\
\midrule
Style 1 & 619 & \tbd & \tbd & \tbd & \tbd & \tbd \\
Style 2 & 381 & \tbd & \tbd & \tbd & \tbd & \tbd \\
Style 3 & 46 & \tbd & \tbd & \tbd & \tbd & \tbd \\
\bottomrule
\end{tabular}%
\missingtable{task4_by_style.tex}{outputs/task4/tables/by_style.tex}}
\end{table}

\begin{table}[t]
\caption{Task 4 test metrics per photograph source and style. The
\texttt{photo3}/Style~1 combination never occurs in the training split
(Section~\ref{sec:data}).}
\label{tab:task4source}
\centering
\resulttable{task4_by_source_style.tex}{%
\begin{tabular}{@{}llccccc@{}}
\toprule
Source & Style & $n$ & $L_1$ & SSIM & PSNR & Edge ratio \\
\midrule
photo3 & Style 1 & 179 & \tbd & \tbd & \tbd & \tbd \\
\bottomrule
\end{tabular}%
\missingtable{task4_by_source_style.tex}{outputs/task4/tables/by_source_style.tex}}
\end{table}

\textbf{What these metrics can and cannot support.} SSIM~\cite{wang2004image} is
paired, interpretable and the metric the FS2K benchmark itself reports, alongside
the sketch-specific SCOOT measure~\cite{fan2019scoot,fan2022facial}; since SCOOT
is not implemented here, the numbers above are \emph{not} directly comparable
with published FS2K benchmark results, and no such comparison is claimed. $L_1$
and PSNR measure pixel fidelity to the artist's drawing, and the edge ratio
measures directly the blur that an $L_1$ term induces, which a pixel metric
rewards rather than penalises. Two perceptual metrics were considered and
deliberately \emph{not} used. FID~\cite{heusel2017gans} is biased at finite sample
sizes, and Chong and Forsyth~\cite{chong2020effectively} show that the bias depends
on the model being evaluated, so one model can beat another purely through a
smaller bias---an effect that dominates at the $n = 46$ of style~3; Parmar et
al.~\cite{parmar2022aliased} add that resizing and compression choices alone
induce large FID variations. KID~\cite{binkowski2018demystifying} would be better
behaved at small $n$, but like FID it requires additional packages and
ImageNet-trained weights downloaded at run time, which the project did not add.
LPIPS~\cite{zhang2018unreasonable} is paired and well defined on 46 images, but its
perceptual calibration comes from patches of natural photographs, not from line
drawings, so its judgement of sketch quality is not established. More generally,
GAN metrics disagree~\cite{borji2019pros,lucic2018gans}, so small differences
should not be over-read, and visual evidence carries more weight here than in
Tasks~1--3.

\begin{figure*}[t]
\centering
\resultfig[0.24\textheight]{task4_style_swap.png}{0.62\textwidth}{From
\path{outputs/task4/figures/style_swap.png}: one test photograph per style,
with the ground truth and the sketch generated in each of the three styles.}
\caption{The same photograph rendered in all three styles, next to the
ground-truth sketch. The style condition is the only input that differs across
the three generated columns.}
\label{fig:task4grid}
\end{figure*}

\begin{figure}[t]
\centering
\resultfig[0.24\textheight]{task4_test_examples.png}{\columnwidth}{From
\path{outputs/task4/figures/test_examples.png}: the two test pairs closest to the
median score of each style.}
\caption{Typical Task 4 test results (median sketch score per style): photograph,
ground truth and generated sketch.}
\label{fig:task4examples}
\end{figure}

On the 1{,}046 official test pairs the generator reaches $L_1$ 0.110, SSIM 0.457
and PSNR 15.33~dB (macro over styles: 0.109, 0.463, 15.44~dB;
Table~\ref{tab:task4}). Style~2, drawn with dense dark strokes, is the hardest
($L_1$ 0.154, SSIM 0.390), styles 1 and 3 are close ($L_1$ 0.085 and 0.087). The
edge ratio is below one for every style (0.56, 0.53, 0.70): the generated strokes
are softer than the artist's, the blur that the $L_1$ term induces and the adversarial
term only partly removes. \textbf{The condition is used.} Generating each test
photograph with a wrong style raises $L_1$ from 0.110 to 0.152, a gap of 0.042 (38\%),
and the style-swap grid (Fig.~\ref{fig:task4grid}) shows three clearly different
renderings of the same photograph---light contours, dark dense strokes, grey
shading. \textbf{The source/style confound does not dominate.} The 180 test pairs
whose (source, style) combination never occurs in training score as well as the
866 seen pairs ($L_1$ 0.110 against 0.110, SSIM 0.467 against 0.455), so style is
not merely inferred from the photograph. Within Style~1 the unseen
\texttt{photo3} pairs have a higher $L_1$ than \texttt{photo1} (0.110 against 0.075,
Table~\ref{tab:task4source}), a difference of the size expected from their
grey-paper background offset of roughly ten grey levels (Section~\ref{sec:data}).

\subsection{Failure analysis}

\begin{figure}[t]
\centering
\resultfig[0.24\textheight]{task4_failure_cases.png}{\columnwidth}{From
\path{outputs/task4/figures/failure_cases.png}: the six test pairs with the lowest
sketch score.}
\caption{Task 4 failure cases (lowest sketch score); automatic notes on each case
are written to \texttt{outputs/task4/eval/failure\_cases.csv}.}
\label{fig:task4failures}
\end{figure}

All six lowest-scoring test pairs (Fig.~\ref{fig:task4failures}) are
\texttt{photo1} portraits with long, voluminous hair, five of them Style~2, and all
carry the automatic note ``strokes too soft'' (edge ratios 0.38--0.58): where the
artist drew dense, dark hair strokes the generator produces thinner grey ones,
the $L_1$-induced blur at its most visible. The worst case is labelled Style~1 but
its ground truth is drawn with heavy dark strokes more typical of Style~2; the
generator follows the label and draws light contours, so part of this error
reflects variation in the annotation rather than in the generator.

\emph{Out-of-domain photographs.} Run through the application on a new photograph
outside FS2K---a short-haired man in front of dense foliage, loosely framed---the
generator captures the facial layout (hair line, eyes, nose, mouth, jaw) but, with
Style~3 selected, adds long strands of hair at the sides of the face. This is the
style/source confound of Section~\ref{sec:fs2kprep} made visible: every Style~3
training sketch comes from the photo2/photo3 sources, mostly long-haired subjects,
so the generator has partly learned ``Style~3'' as that appearance rather than as a
drawing technique. Styles~1 and~2, trained on the varied CASIA-WebFace portraits,
are the better choice for such inputs. Busy backgrounds and loose framing add a
second domain gap, since FS2K photographs are tight crops on plain backgrounds; a
face-detection crop before inference and training data that crosses styles with
sources would address both.

\section{Experimental Setup}
\label{sec:setup}

\textbf{Hardware.} All models are trained with PyTorch~\cite{paszke2019pytorch} on
NVIDIA T4 GPUs~\cite{nvidiat4} with mixed precision~\cite{micikevicius2018mixed}
(fp16 autocast with gradient scaling, since the T4 has no bf16 support), in two
environments. Task~1 (search, final training, evaluation and export) and both
Task~2 searches ran on a free-tier Google Colab~\cite{colab} instance with one
T4. On the deadline day the free Colab GPU quota was exhausted, so the remaining
work was moved to a Kaggle notebook with two T4 GPUs. A first Kaggle run produced
working models for every task, but inspection showed over-smoothed restorations and
faint sketches, so \emph{all final models were retrained} in a second Kaggle run,
which is the only one reported. A single runner, \texttt{scripts/kaggle\_run.py},
prepares both datasets and drives three jobs in parallel: GPU~0 trains and
evaluates the three Task~1 variants and writes the skip-ablation table; GPU~1 runs
Task~4 and, alongside it, Task~2 (classifier and the three skip specialists)
followed by Task~3, which is built from the new Task~2 models. Every step writes a
completion marker, so the runner can simply be restarted after an interruption.
The final models use the hyperparameters selected by the Optuna studies (the
Task~1 and Task~2 studies ran on Colab and are committed as
\texttt{configs/task1.yaml} and \texttt{configs/task2\_*.yaml}); the deadline
forced a shorter Optuna budget for Task~3 (Table~\ref{tab:budgets}).

\textbf{Resumable runs.} Because free-tier sessions
disconnect without warning, every component is built to resume: training writes a
complete checkpoint every epoch (model, optimiser, scheduler, gradient scaler,
early-stopping state, history and the Weights \& Biases run id) and
\texttt{--resume} continues from it with the stored configuration treated as
authoritative, so a resumed run cannot silently change its schedule. The best
checkpoint is written before the last one, so a crash between the two can never
leave the best behind. Data order and runtime corruptions depend only on the seed
and the epoch index, so a resumed run sees the same data as an uninterrupted one;
a test interrupts a run mid-epoch and asserts that the resumed weights match the
uninterrupted ones exactly.

Hyperparameter optimisation uses Optuna~\cite{akiba2019optuna,optunadocs}. Its
define-by-run interface allows the search space to be built dynamically, which is
what lets conditional parameters (such as a latent shape that implies a depth) be
expressed directly. All studies use the TPE
sampler~\cite{bergstra2011algorithms,watanabe2023tree} seeded at 42, which fits
one density to the parameters of the best trials and another to the rest and
maximises their ratio; random search would be the weaker
baseline~\cite{bergstra2012random}. Pruning uses the median stopping
rule~\cite{golovin2017google}, which stops a trial whose best intermediate value
is worse than the median of completed trials at the same step and which Vizier
reports gives a two- to three-fold speedup while still finding the best trial.
Successive halving and Hyperband~\cite{jamieson2016non,li2018hyperband,li2020system}
were considered and not used: with 20--40 trials on one T4, their aggressive
early stopping would spend most of the budget on very short runs whose ranking is
unreliable for GAN and mixture training, whereas the median rule is model-free
and matches the fixed-epoch validation schedule. Task~3 adds the task-specific
routing-collapse criterion of Section~\ref{sec:task3}, which no generic pruner
can express. Each study's SQLite database lives on local disk and a consistent
snapshot is copied to persistent storage (Google Drive on Colab, the notebook's
working directory on Kaggle) after every finished trial, so a search survives a
disconnect and only the remaining trials are run. Table~\ref{tab:budgets} lists
what each study actually ran.

\begin{table}[t]
\caption{Optuna budgets and final training runs. C/P: completed and pruned
trials; time: sum of trial durations. The Task~3 budget was cut on Kaggle; the
original plan was 20 trials of 1+4 epochs and a 2+20-epoch final run.}
\label{tab:budgets}
\centering
\footnotesize
\setlength{\tabcolsep}{2pt}
\begin{tabular}{@{}lllcll@{}}
\toprule
Study & GPU & Trials $\times$ ep. & C/P & Time & Final run \\
\midrule
T1 UDAE & Colab & 33 (of 40) $\times$ 15 & 15/18 & 89.6 min & $3{\times}60$ ep., Kaggle \\
T2 classifier & Colab & 30 $\times$ 12 & 5/25 & 15.6 min & 60 ep., Kaggle \\
T2 specialists & Colab & 20 $\times$ 6 ($\times 3$) & 12/8 & 61.5 min & $3{\times}40$ ep., Kaggle \\
T3 soft MoE & Kaggle & 6 $\times$ (1+3) & 5/1 & 19.3 min & 2+10 ep. \\
T4 cGAN & Kaggle & 18 $\times$ 12 & 10/8 & 31.0 min & 200 ep. \\
\bottomrule
\end{tabular}
\\[2pt]

\end{table}

\textbf{Experiment tracking} uses Weights \& Biases~\cite{wandb}, in the project
\texttt{genai-a1}
(\url{https://wandb.ai/mustafafaraz87-fast-nuces/genai-a1}). Every Optuna trial is
its own run, grouped per study, with its parameters in the configuration and its
per-epoch validation metrics logged; final runs additionally log per-condition
metrics, sample image grids at fixed intervals using \emph{fixed} validation
indices so that progress is comparable across epochs, and the best checkpoint as
a model artefact. The Colab and Kaggle runs log to the same project.


\section{Application Architecture}
\label{sec:app}

The four tasks are served by one application, \emph{GenAI Studio}, with four
workspaces named as in the brief: \textbf{Universal Restoration} (Task~1),
\textbf{Hard-Routed Restoration} (Task~2), \textbf{Soft Mixture-of-Experts
Restoration} (Task~3) and \textbf{Face-to-Sketch Generator} (Task~4).
Figure~\ref{fig:architecture} shows its structure.

\subsection{Interface design}

The interface was designed in Google Stitch~\cite{stitch} before any frontend code
was written, as the brief requires; Fig.~\ref{fig:stitch} reproduces the original
Stitch screens, and the exported HTML is kept in \texttt{design/stitch/}. The
design defines a four-tab shell over a 12-column grid: a narrow input card (upload
dropzone, sample images, \emph{Corruption} and \emph{Severity} segmented controls
and the action button) and a wide results column with \emph{Original},
\emph{Model input} and \emph{Restored output} panels, a details card with a
download control and an error map. Each workspace adds what its task exposes: a
\emph{Predicted}/\emph{Oracle} routing toggle and classifier probability bars for
hard routing; four expert weights, a ``main contributor'' and the four branch
outputs for the mixture; and a webcam button and a \emph{Style 1/2/3} selector
for the sketch generator. The implementation keeps the design's Tailwind
configuration (colour roles, radii, spacing and type scale) unchanged, but every
number in the mock-ups---GPU names, latencies, LPIPS scores---was a Stitch
placeholder and was replaced by values returned by the API or removed: the
application shows only real data.

\begin{figure*}[t]
\centering
\resultfig{stitch/universal.png}{0.40\textwidth}{Copied from
\path{design/stitch/universal.png}.}\hspace{0.03\textwidth}%
\resultfig{stitch/hard-routed.png}{0.40\textwidth}{Copied from
\path{design/stitch/hard-routed.png}.}\\[4pt]
\resultfig{stitch/moe.png}{0.40\textwidth}{Copied from
\path{design/stitch/moe.png}.}\hspace{0.03\textwidth}%
\resultfig{stitch/sketch.png}{0.40\textwidth}{Copied from
\path{design/stitch/sketch.png}.}
\caption{The original Google Stitch designs of the four workspaces, as generated
in Stitch before implementation (evidence required by the brief): Universal
Restoration (top left), Hard-Routed Restoration (top right), Soft
Mixture-of-Experts Restoration (bottom left) and Face-to-Sketch Generator (bottom
right). The numbers and model names in these mock-ups are Stitch placeholders;
the implemented application replaces them with values returned by the backend.}
\label{fig:stitch}
\end{figure*}

\begin{figure}[t]
\centering
\begin{tikzpicture}[
  font=\scriptsize,
  box/.style={draw, rounded corners=2pt, align=center, inner sep=3pt,
    text width=0.80\columnwidth},
  arr/.style={-{Stealth[length=4pt]}, semithick},
  node distance=4.5mm]
\node[box] (browser) {\textbf{Browser}: React + TypeScript + Tailwind CSS\\
four workspaces, webcam capture, error map, downloads};
\node[box, below=of browser] (nginx) {\textbf{frontend container} (nginx, port 5173)\\
built static files; \texttt{/api/*} proxied to the backend (same origin)};
\node[box, below=of nginx] (api) {\textbf{backend container} (FastAPI, port 8000)\\
upload validation $\cdot$ preprocessing and corruption (\texttt{src/data}, as in
training) $\cdot$ ONNX Runtime sessions (CPU) $\cdot$ PSNR/SSIM and timing};
\node[box, below=of api] (models) {\textbf{\texttt{./models}} mounted read-only at
\texttt{/models}: \texttt{udae}, \texttt{classifier},
\texttt{specialist\_\{salt,blur,occlusion\}}, \texttt{moe}, \texttt{generator}
(\texttt{.onnx} + JSON model cards)};
\node[box, below=of models] (train) {\textbf{Training} (Colab / Kaggle T4): PyTorch,
Optuna, Weights \& Biases $\rightarrow$ ONNX export with parity check};
\draw[arr] (browser) -- node[right, font=\tiny] {HTTP} (nginx);
\draw[arr] (nginx) -- node[right, font=\tiny] {\texttt{/api}} (api);
\draw[arr] (models) -- node[right, font=\tiny] {loaded at start-up} (api);
\draw[arr] (train) -- node[right, font=\tiny] {exported files} (models);
\node[draw, dashed, rounded corners=3pt, inner sep=2.5pt, fit=(nginx) (api) (models),
  label={[font=\tiny, anchor=south east]north east:\texttt{docker compose up --build}}] {};
\end{tikzpicture}
\caption{Application architecture. One command starts both containers; the
backend reuses the training code's corruption and preprocessing functions, so the
application degrades and prepares images exactly as training did.}
\label{fig:architecture}
\end{figure}

\subsection{Backend}

The backend is a FastAPI application~\cite{fastapi} that runs every model through
ONNX Runtime~\cite{onnxruntime} on the CPU execution provider, so it needs no GPU
and does not import PyTorch at all (a test asserts it). Its endpoints are listed in
Table~\ref{tab:endpoints}. The health endpoint reports status, uptime, runtime
versions, every model with its model card, and which workspaces are ready; the
options endpoint returns the corruption types, fixed levels, permitted custom
parameter ranges and style names, so the interface builds its controls from the
server rather than duplicating constants. Crucially, the backend imports
\emph{the same} corruption and preprocessing code as training
(\texttt{src/data/}), so an image passed through the application is degraded and
prepared exactly as it would have been during training, and the corrupted input
is rounded to 8 bits so that the image returned as ``Model input'' is exactly
what the model saw.

\begin{table}[t]
\caption{Backend API (all paths under \texttt{/api}; every request is multipart
form data and every image in a response a PNG data URL).}
\label{tab:endpoints}
\centering
\footnotesize
\setlength{\tabcolsep}{3pt}
\begin{tabular}{@{}l>{\raggedright\arraybackslash}p{4.55cm}@{}}
\toprule
Endpoint & Purpose \\
\midrule
\texttt{GET /health} & status, uptime, runtime versions, models and model cards, ready workspaces \\
\texttt{GET /options} & corruption types, levels, training ranges, styles, upload limits \\
\texttt{GET /samples} & bundled pet photographs from the official test split \\
\texttt{POST /corrupt} & runtime corruption preview (type, level or custom parameters, seed) \\
\texttt{POST /restore/universal} & Task 1, \texttt{udae.onnx} \\
\texttt{POST /restore/hard} & Task 2, classifier + one specialist or identity; predicted or oracle routing \\
\texttt{POST /restore/moe} & Task 3, \texttt{moe.onnx}: output, four weights, four branch outputs \\
\texttt{POST /sketch} & Task 4, \texttt{generator.onnx}, style 1, 2 or 3 \\
\bottomrule
\end{tabular}
\end{table}

Uploads are validated by decoding rather than by declared content type: JPEG,
PNG, WEBP and BMP are accepted up to 10~MB and 50 megapixels, which rejects
decompression bombs; EXIF rotation is applied and greyscale, palette and
transparent images are converted to RGB. Parameters outside the training ranges
are refused with an explanatory message rather than silently clamped, so the
evaluator cannot unknowingly test the models outside the distribution they were
trained on. Requests may supply either an uploaded file or a bundled sample, and
either a fixed severity level or custom parameters, never both. When the server
applies the corruption itself, the response also contains PSNR and SSIM of the
model input and of the output against the clean original, with SSIM implemented
in NumPy to the training definition (agreement with the PyTorch version within
$10^{-5}$ is tested), and every response reports the inference and total server
time---for hard routing split into classifier and expert time. A missing model
file disables only the requests that need it, with a message naming the file.

\subsection{Frontend and deployment}

The frontend is a React~\cite{react} and TypeScript application styled with
Tailwind CSS~\cite{tailwind} and built with Vite. It implements the four Stitch
workspaces against the API, each at its own URL and kept mounted when hidden, so
inputs and results survive switching tabs. Each restoration workspace accepts an
upload or a sample and offers runtime corruption with a fixed level or custom
parameters and a reproducible seed, a corruption preview, PSNR and SSIM under the
image panels, an absolute error map computed in the browser and a download of the
result; the hard-routing workspace adds the classifier probabilities, the
predicted and true class, the expert used and the classifier/expert timing, and
the mixture workspace the four expert weights and branch outputs. The
Face-to-Sketch Generator accepts an upload or a webcam capture and a style, shows
photograph and sketch side by side and offers the sketch for download. Fonts and icons
are self-hosted from npm packages and Tailwind is compiled at build time, so the
application needs no external CDN.

Frontend and backend each run in their own container and start together with one
documented command, \texttt{docker compose up --build}~\cite{dockercompose}. The
application is then served at \url{http://localhost:5173}; the frontend container
(nginx) serves the built files and forwards \texttt{/api} to the backend
container, so the browser talks to a single origin and no CORS configuration is
needed. The backend publishes port 8000 with interactive API documentation, runs
as a non-root user with a health check, and the frontend starts once that check
passes. The trained ONNX models are not baked into the image: the
\texttt{./models} folder is mounted read-only, so updated models need only a
restart, and without them the backend still starts, reports itself as degraded
and the interface names the missing files.

\begin{figure*}[t]
\centering
\resultfig[0.18\textheight]{app_universal.png}{0.40\textwidth}{Screenshot of the
running application with the trained models: Universal Restoration.}\hspace{0.03\textwidth}%
\resultfig[0.18\textheight]{app_hard.png}{0.40\textwidth}{Screenshot:
Hard-Routed Restoration, with classifier probabilities.}\\[4pt]
\resultfig[0.18\textheight]{app_moe.png}{0.40\textwidth}{Screenshot: Soft
Mixture-of-Experts Restoration, with the expert weights.}\hspace{0.03\textwidth}%
\resultfig[0.18\textheight]{app_sketch.png}{0.40\textwidth}{Screenshot:
Face-to-Sketch Generator, photograph and sketch side by side.}
\caption{The four workspaces of the running application with the trained models.}
\label{fig:screenshots}
\end{figure*}

\subsection{ONNX export and parity}

Seven ONNX~\cite{onnx,pytorchonnx} files are exported, at opset 17 in float32 with
a dynamic batch axis: the universal autoencoder, the classifier, the three
specialists, the complete mixture pipeline as a single graph (gate, temperature
constant, softmax, identity branch, three experts and the weighted sum, emitting
the output, the four weights and all four branch outputs) and the sketch
generator; the discriminator is a training component only. The TorchScript-based
exporter is used explicitly, because the newer default exporter of recent PyTorch
versions targets opset 18 and can emit invalid \texttt{Split} nodes. Every export is verified against its PyTorch original in
evaluation mode on at least 32 \emph{real} test inputs spanning all conditions and
severities---not random noise, so that the gate sees realistic logits---and the
maximum absolute difference must stay below $10^{-4}$, or the file is rejected
and the script fails. The measured maximum differences are $7.2\times10^{-7}$ (universal
autoencoder), $1.2\times10^{-7}$ (classifier probabilities), $1.6\times10^{-6}$,
$2.1\times10^{-6}$ and $9.4\times10^{-6}$ (salt-and-pepper, blur and occlusion
specialists), $1.1\times10^{-5}$ (mixture output; $1.2\times10^{-7}$ for its weights)
and $2.5\times10^{-6}$ (generator), all at least nine times inside the tolerance. Each model
is accompanied by a JSON model card recording its configuration, input and output
specification, preprocessing, test metrics, hyperparameters, the measured
parity, the Weights \& Biases run URL, the git commit and the export time; the
backend shows these cards in the application's system-information panel.


\section{Limitations}
\label{sec:limitations}

\emph{Severity coverage.} All three test ``high'' levels sit at the upper edge of
the training ranges, and no training blur ever reaches the strength of the
$(7, 2.5)$ test configuration (Section~\ref{sec:data}). Results at the high
level therefore measure extrapolation to the boundary of the training
distribution as much as restoration ability; for Task~1 this shows as a 0.8~dB
drop from medium to high blur and a steady decline for occlusion, while
salt-and-pepper is unaffected.

\emph{Search fidelity and short-trial bias.} Hyperparameters were selected from
short trials---15 epochs for Task~1, 12 for the Task~2 classifier, 6 per
specialist, 4 for Task~3 and 8 for Task~4---and applied to runs three to fifteen
times longer. Short budgets favour configurations that learn quickly, and the
evidence is visible: both restoration searches chose the smallest batch size,
which gives the most optimiser steps per epoch, and Task~1's selected
configuration gained 0.076 in validation score between its 15-epoch trial and
the 60-epoch run. The selected configurations are good regions of the search
space found under a budget, not proven optima; Task~1's best trial was also its
last, so more trials could still have helped.

\emph{Reduced budgets for Tasks 3 and 4.} The free Colab GPU quota ran out on the
deadline day, and Tasks~3 and~4 were trained on Kaggle with budgets cut to fit:
six Optuna trials of one warm-up and three joint epochs and a final run of
2+10 epochs for Task~3 (planned: 20 trials, 2+20 epochs), and ten trials of eight
epochs for Task~4 (planned: 20 of ten). With ten random start-up trials
configured for TPE, neither study reached the model-based phase of the sampler,
so both are small random searches around the brief's reference configuration.
Their hyperparameters are therefore less well optimised than those of Tasks~1
and~2, and Task~3's shorter joint fine-tuning may understate what soft routing
can achieve.

\emph{The Task~1 versus Task~2 bottleneck confound.} The specialists' search
offered $16\times16$ latents and selected one of 8{,}192 values, twice the 4{,}096
of the universal model, whose range was capped to keep a meaningful bottleneck.
Part of any advantage of hard routing over the universal model may therefore come
from the larger latent rather than from specialisation; a controlled comparison
would retrain the universal model with the same latent shape.

\IfFileExists{tables/task1_skip_ablation.tex}{}{\emph{Missing ablation.} The
limited-skip ablation of Task~1 is implemented but was not trained within the GPU
budget, so the leak-versus-detail trade-off of skip connections is argued and
bounded by the two diagnostics rather than measured.}

\emph{Mixed precision for the GAN.} Task~4 trains with fp16 autocast and separate
gradient scalers for the generator and the discriminator. Its numerical
stability could not be verified before the Kaggle run, since the development
machine had no GPU; a non-finite loss aborts the run (or prunes the trial) rather
than producing silent NaNs, and full precision remains available as a fallback.


\emph{Scale.} With 2{,}944 training images at $128 \times 128$ and models of a
few million parameters, none of these systems is comparable with published
all-in-one restoration models~\cite{li2022allinone,potlapalli2023promptir}, which
are trained on orders of magnitude more data. No such comparison is made.

\emph{The mixture objective.} As Section~\ref{sec:task3} notes, training the
blend on its reconstruction error is the cooperative formulation that Jacobs et
al.\ identify as discouraging specialisation~\cite{jacobs1991adaptive}. The
pre-trained experts and the cross-entropy term counteract it, but the system is
not an independent demonstration that soft routing induces specialisation; the
specialisation was supplied by Task~2.

\emph{Balance regulariser blind spot.} Equation~(\ref{eq:balance}) sees only
batch-mean weights and cannot detect a gate that routes each class confidently to
the wrong branch. The routing matrix and the cross-entropy term cover that case,
but the regulariser alone does not.

\emph{FS2K confounds and sample sizes.} Style is confounded with photograph
source in the official FS2K training split---styles 1 and 2 come only from
\texttt{photo1}, style~3 only from \texttt{photo2} and \texttt{photo3}---so part
of what looks like style control may be source recognition; the test set
contains only 46 style-3 pairs, which makes distributional metrics unusable for
that style and widens the error bars on every per-style number; and the 179
\texttt{photo3} style-1 test pairs carry the grey-paper background offset
(Section~\ref{sec:data}).

\emph{Metrics.} PSNR correlates with subjective quality only within a fixed
content set~\cite{huynh2008scope}, and SSIM and PSNR are shallow
functions~\cite{zhang2018unreasonable}. Task~4 reports no FID and no LPIPS: FID is
biased in a model-dependent way at small sample sizes~\cite{chong2020effectively}
and would be meaningless on the 46 style-3 pairs, LPIPS is calibrated on natural
photographs rather than line drawings, and both need extra packages and
downloaded ImageNet weights that the project did not add. Sketch quality is
therefore judged by $L_1$, PSNR, SSIM, the edge ratio and visual evidence, which
measure fidelity to one artist's drawing rather than perceptual quality.
Conclusions are drawn from several metrics together with visual evidence, never
from a single number.

\emph{Deployment.} ONNX inference runs on CPU in the application container, so
reported inference times reflect the evaluation machine and are not comparable
with GPU timings.

\section{Conclusion}

Three restoration designs were built on one data pipeline and compared on
identical test entries. Explicit routing justified its complexity: a classifier
that is right on 99.8\% of inputs lets three specialists beat the universal model
by 1.2~dB on the corrupted entries and return 99\% of clean images exactly, which a
universal autoencoder can never do. Making the routing soft and training it
jointly gained a further 0.33~dB, consistently across corruptions, but not through
blending: the learned gate is nearly one-hot, the temperature is compensated by
the logit scale, and on images with two corruptions the gate still picks one
expert. The gain comes from fine-tuning gate and experts together. Of the
research-driven decisions, the experiments confirmed the SSIM-heavy loss weighting
(both searches moved $\alpha$ far below the brief's 0.8, consistent with Zhao et
al.), resize-convolution (no checkerboard artefacts appear), the balanced sampler
with the brief's balance term (no collapse, usage within 0.004 of $1/4$) and decoder
FiLM with a projection discriminator (a wrong style raises $L_1$ by 38\%, and
unseen source/style combinations generalise). One decision was revised by
evidence: the pure bottleneck. The ablation showed that one limited skip at
$32\times32$ restores detail (detail ratio 0.48 to 0.81, +4.5~dB on the corrupted
entries) while impulse survival falls, so the deployed models keep the $8\times8$
latent but add that skip. The largest remaining weaknesses are occlusion, where
every system is limited by content that must be invented, and the soft strokes of
the sketch generator, which a perceptual or sharper adversarial objective would be
the next step to address.

\appendix[Use of AI assistance]

AI tools were used extensively during this work, as the brief permits, and
Table~\ref{tab:aiuse} lists them. The main one was Claude Code (Anthropic, with
Claude Opus models), used as a coordinating coding assistant that split the work
into parallel sub-agents, each with a narrow assignment and its own tests: one
for research and citation verification, one for the data pipeline and the shared
training utilities, one for each task's training, evaluation and export code, and
one each for the FastAPI backend, the React frontend, the Docker setup and the
drafting of this report. The interface was designed with Google Stitch.

Nothing generated was accepted on trust. Every component ships with unit and
smoke tests---111 Python tests for the data pipeline, models and the four tasks,
97 backend tests and 54 frontend tests---and the coordinating session re-ran every
sub-agent's tests independently before anything was committed. Every ONNX export
must pass its parity check. Every one of the 114 references was checked against
Crossref, arXiv or the publisher's or proceedings page, with the verifying URL
recorded next to its BibTeX entry, and quoted claims were spot-checked
independently. Verification found and corrected real errors, five of which
illustrate the process:

\begin{enumerate}
\item \emph{Blur severity mislabelled.} Blur severity was first assigned from
$\sigma$ alone. Measuring the strength of the kernel actually applied
(Eq.~\ref{eq:blurstrength}) showed that a wide Gaussian truncated to three taps is
milder than the medium test level although its $\sigma$ is larger; relabelling by
measured strength corrected 86 of the 184 validation blur labels.
\item \emph{SSIM convention.} A cross-check against \texttt{torchmetrics} found a
small disagreement (about $3\times10^{-5}$). The cause is a convention, not a bug:
\texttt{torchmetrics} pads the image and averages over all windows, whereas the
implementation here follows Wang et al.'s reference and averages over the valid
region only. The difference is documented, and the test now asserts agreement
with \texttt{torchmetrics}' SSIM map cropped to the valid region.
\item \emph{SQLite connection leak on Windows.} In the code that snapshots each
Optuna study to persistent storage, a \texttt{sqlite3} connection used as a
context manager commits but is not closed, which on Windows kept the snapshot file
locked; it is now closed explicitly.
\item \emph{A quick check that contaminated real checkpoints.} The first Task~2
classifier evaluation reported test accuracy 0.300, macro-recall exactly 0.250 and
validation macro-F1 0.1000---one class predicted for every input. Checking these
numbers against what a trained model must produce exposed the cause: a
\texttt{--smoke} run had saved a finished tiny model into the real checkpoint
folder, so the real run, started with \texttt{--resume}, found ``nothing to do'' and
evaluation silently used the tiny model. Real runs now ignore smoke checkpoints,
evaluation and export refuse them, and a regression test covers the case. This is
the clearest example of generated code that passed its own tests yet failed in
the real workflow.
\item \emph{Zhao et al.'s $\alpha$.} Their recommended $\alpha = 0.84$ looks like
support for the brief's $\alpha = 0.8$, but checking the paper showed that it
weights the MS-SSIM term, not $L_1$, so it corresponds to $\alpha \approx 0.16$
here. The search range was set to contain both values, and the selected 0.285 is
interpreted with the correct parameterisation (Section~\ref{sec:task1}).
\end{enumerate}

The author reviewed all submitted code, text and figures, ran and checked the
experiments, and is responsible for the correctness of the complete submission.

\begin{table*}[t]
\caption{Use of AI assistance: tools, what they were used for, and how their
output was tested or corrected.}
\label{tab:aiuse}
\centering
\footnotesize
\begin{tabular}{@{}>{\raggedright\arraybackslash}p{2.7cm}%
>{\raggedright\arraybackslash}p{5.6cm}>{\raggedright\arraybackslash}p{8.6cm}@{}}
\toprule
Tool & Used for & How the output was tested or corrected \\
\midrule
Claude Code (Anthropic), coordinating session & Planning, splitting the work
into parallel sub-agents, integration, debugging, commits & Each sub-agent's tests
re-run independently before every commit; the student reviewed diffs, ran the
training on Colab and Kaggle and checked results against expectations
(correction 4) \\
\addlinespace[2pt]
Claude Code sub-agent: research & Literature search, comparison of design
alternatives, collection of the 114 references & Every citation checked against
Crossref, arXiv or the proceedings page (URL kept in \texttt{references.bib});
quotes spot-checked independently; parameterisation errors corrected (5) \\
\addlinespace[2pt]
Claude Code sub-agents: data and shared code & Oxford-IIIT Pet split, runtime
corruptions, manifests, balanced sampler, FS2K preparation; losses, metrics,
checkpoints, Optuna and W\&B helpers, ONNX export & 44 tests (corruptions,
FS2K, models, shared utilities), e.g.\ blur equal to \texttt{torchvision} within
$10^{-5}$, exact resume after interruption; corrections 1--3 \\
\addlinespace[2pt]
Claude Code sub-agents: Tasks 1--4 & Models, training, Optuna studies,
evaluation, figures and ONNX export for each task & 67 CPU smoke and unit tests
(14, 14, 17 and 22 per task), ONNX parity at most $10^{-4}$ on real inputs,
evaluation numbers sanity-checked (correction 4) \\
\addlinespace[2pt]
Claude Code sub-agents: backend, frontend, Docker & FastAPI service, React +
Tailwind interface implementing the Stitch design, Dockerfiles and Compose file
& 97 backend tests (validation, errors, endpoints, metrics, missing models), 54
Vitest tests, end-to-end run under Docker Compose with placeholder models, and the
exported Task~1 model checked through the real backend \\
\addlinespace[2pt]
Claude Code sub-agent: report & Drafting this report from the design notes and
the experiment log & Every number traced to the experiment log or to the output
files of the scripts; text reviewed by the student \\
\addlinespace[2pt]
Google Stitch & Visual design of the four workspaces
(Fig.~\ref{fig:stitch}) & Design reviewed; its Tailwind configuration reused;
placeholder numbers replaced by data from the API \\
\bottomrule
\end{tabular}
\end{table*}

\section*{Availability}

The complete implementation---data preparation, training, Optuna search,
evaluation and ONNX export scripts, backend, frontend, Dockerfiles and Docker
Compose configuration---is available at \url{\repourl}; the application starts
with \texttt{docker compose up --build} once the ONNX models are placed in
\texttt{models/} (\url{https://github.com/fsdev87/generative-ai-A1/releases/latest/download/models.zip}). Experiment
tracking records are in the Weights \& Biases project
\url{https://wandb.ai/mustafafaraz87-fast-nuces/genai-a1}. The demonstration video
is at \todo{YouTube link}.

\bibliographystyle{IEEEtran}
\bibliography{references}

\end{document}
